\chapter{可泛化前馈重建}

\paragraph*{本章阅读锚点}
原始 3DGS 是逐场景优化方法：给定一组已知位姿图像，经过数千到数万步梯度下降得到一个场景的高斯集合。可泛化前馈重建要改掉这个前提：模型在大数据上预训练，测试时输入少量图像，一次或少数几次前向传播直接输出 3D Gaussians。早期方法多采用像素对齐预测，即每个输入像素沿相机射线预测一个高斯；随后深度、MVS、体素、锚点、token、global scene token 和 LRM 路线不断出现，目标都是减少像素绑定带来的冗余、漂浮和跨视角不一致。本章按输入假设和表示演进阅读：先看多视角 pose-given 前馈重建，再看 pose-free/uncalibrated，最后看大重建模型、单图和领域化方法。

\section{多视角前馈重建}

多视角前馈重建的基本锚点是 MVS：先估计深度或点图，再把像素反投影为高斯中心，并预测尺度、不透明度和颜色。MVSGaussian、FreeSplat、TranSplat、DepthSplat 等方法都在这条线上；后续 VolSplat、AnchorSplat、TokenGS、GlobalSplat 则试图摆脱像素网格，让高斯数量和位置更贴合 3D 场景本身。



% METHOD: 2405.12218
\paragraph{MVSGaussian（MVS 到高斯的早期前馈锚点）（2024.05）}
MVSGaussian~\cite{2405_12218} 是早期可泛化 3DGS 代表。它用多视角立体视觉聚合跨视角特征，在 3D 点上编码几何感知特征，再由解码器输出高斯位置偏移、颜色、不透明度、尺度和旋转。

它还引入混合渲染思路：利用体渲染的深度一致性增强几何，同时保持 splatting 的速度。多个视角生成的点云可通过几何一致性聚合，作为快速逐场景微调初始化。相对后续 DepthSplat，它的深度基础模型成分较弱，但奠定了 MVS-to-Gaussian 结构。
\textbf{实现细节：} 相对原始 3DGS 的逐场景梯度优化，MVSGaussian 把流程改成前馈预测：训练一个通用网络，推理时从多视角输入图像直接输出场景的 3D 高斯参数。输入包括源视角图像、相机参数和跨视角特征；输出包括像素或点对齐的高斯位置偏移、颜色、不透明度、尺度和旋转。它的关键是把 MVS 代价体深度估计接到高斯初始化上，让每个像素的深度候选变成三维点，再由解码器预测可渲染高斯。
\textbf{实现补充：} 具体管线可以拆成四步：先用 FPN 提取多尺度图像特征并构建跨视角代价体；再估计深度，把像素反投影到 3D 空间并形成像素对齐高斯；随后用混合渲染结合体渲染式深度一致性和 splatting 速度，让训练损失同时约束图像颜色和几何；最后对多视角生成的高斯做一致性聚合，并可作为快速逐场景微调的初始化。它的边界在于前馈模型依赖训练数据分布和相机标定质量，复杂遮挡或宽基线场景仍可能需要后续优化。
\textbf{方法信息：} 引用很多，Semantic Scholar 约 108 次。
\textbf{设计动机与边界：} 相对原始 3DGS 以独立高斯和图像重建误差为主的设定，MVSGaussian 把问题落在以下可实现的改动上：MVSGaussian的核心设计包含三个部分：(1) MVS-based高斯编码：利用多视角立体视觉的几何感知能力，通过特征金字塔网络提取多尺度特征，在3D点上聚合跨视角特征，编码几何感知的高斯表示，然后通过解码器将特征解码为高斯参数（位置偏移、颜色、不透明度、缩放、旋转等），以前馈方式直接回归高斯参数；(2) 混合高斯渲染：为提升性能，将高效的体渲染设计与高斯splatting相结合——利用体渲染的深度一致性约束增强几何质量，同时保持splatting的渲染速度优势；(3) 多视角几何一致聚合：为支持快速场景特定微调，将可泛化模型生成的多个视角的点云通过几何一致性检查进行有效聚合，作为per-scene优化的初始化，显著减少优化所需的迭代次数。相比可泛化NeRF方法需要分钟级微调和秒级渲染，MVSGaussian实现实时渲染和更优合成质量。其设计目的不是单独增加一个模块，而是用上述表示、监督或渲染约束消除原流程在目标任务中的具体失配。提出首个基于MVS的可泛化3DGS方法，通过混合渲染设计结合splatting和体渲染优势，支持实时渲染、高质量合成和快速场景特定优化，在四个基准数据集上达到SOTA适用时还应保留论文给出的边界：MVS特征提取增加了前馈推理的计算开销；对输入视角数量有一定要求（至少3视角）；在纹理缺失区域的几何质量受限于MVS能力
% METHOD: 2405.17958
\paragraph{FreeSplat（长序列室内自由视角前馈 3DGS）（2024.05）}
\textbf{锚点细化：} 前述方法（DepthSplat、TranSplat、HiSplat）均假设少量输入视角（通常 2--3 张），生成的高斯数量有限。当输入视角数量增加时，像素对齐高斯在重叠区域产生大量冗余。FreeSplat（2024.05）~\cite{2405_17958}针对室内大场景的长序列重建需求，提出了低成本跨视角聚合框架和像素级三元融合方法。其核心挑战在于：随着输入视角数量增加，像素对齐高斯在重叠区域产生大量冗余，需要一种高效的增量式融合策略来消除重复基元。
FreeSplat~\cite{2405_17958} 针对的是另一种输入变化：不是两三张稀疏图像，而是室内长序列。pixelSplat、MVSplat 等前馈模型在少量视图上可行，但视角一多，每个视角都会在重叠表面预测自己的像素对齐高斯，同一墙面、桌面或地面会被重复生成，计算和存储都浪费。FreeSplat 的核心改变是把长序列聚合和冗余高斯融合放进前馈管线。

它先用低成本 CNN 骨干替代重型跨视角 Transformer，在当前视角和邻近视角之间自适应构造代价体，利用平面扫描立体匹配编码相机姿态信息，并由多尺度编码器-解码器预测深度候选和局部高斯三元组。随后 Pixel-wise Triplet Fusion 将局部三元组与全局三元组对齐，通过深度一致性检查和特征融合消除重叠区域的重复高斯。训练时，FreeSplat 随机改变输入视角数量和视角跨度，并监督更大范围目标视角，使模型不是只学窄基线插值，而是学会自由视角外推。

相对 DepthSplat 的 ``深度先验 + 前馈高斯''，FreeSplat 的锚点价值在于把可泛化 3DGS 推到长序列室内场景，并给出了从局部高斯到全局高斯的融合机制。PTF 不是简单平均重复高斯，而是在潜在特征空间累积多视角观测，使同一三维区域的颜色、几何和置信度一起融合。论文报告冗余高斯减少约 30--50\%，同时在室内渲染质量和深度精度上优于 pixelSplat、MVSplat 等基线。\textbf{方法信息：} 代码：\url{https://github.com/wangys16/freesplat}；仓库提供安装、数据准备和训练/评测说明；GitHub 约 164 stars，Semantic Scholar 约 94 次引用。

\textbf{实现细节：} 前述方法（DepthSplat、TranSplat、HiSplat）均假设少量输入视角（通常 2--3 张），生成的高斯数量有限。当输入视角数量增加时，像素对齐高斯在重叠区域产生大量冗余。FreeSplat（2024.05） 针对室内大场景的长序列重建需求，提出了低成本跨视角聚合框架和像素级三元融合方法。其核心挑战在于：随着输入视角数量增加，像素对齐高斯在重叠区域产生大量冗余，需要一种高效的增量式融合策略来消除重复基元。与 DepthSplat 的架构差异。DepthSplat 使用 Transformer 架构处理固定数量的输入视角，计算量随视角数二次增长。FreeSplat 选择纯 CNN 骨干（EfficientNet/ResNet）替代 Transformer 进行二维特征提取，使计算开销随视角数线性增长，从而支持任意长度的输入序列。在邻近视角间自适应构建代价体（$K{=}128$ 个深度平面，范围 $[d_{\mathrm{near}}{=}0.5, d_{\mathrm{far}}{=}15.0]$ 米），而非全局构建。具体做法是根据位姿邻近度自适应选择 $N$ 个邻近视角，通过平面扫掠和余弦相似度构建代价体，再经 UNet++ 解码器预测深度图。像素级三元融合。每个像素被解投影为三维高斯三元组 $(\mu_l^t, w_l^t, f_l^t)$——中心坐标、融合权重（$(0,1)$ 之间的标量）和潜在特征向量。（$t{=}1, 2, \ldots, T$）：每步先将全局高斯中心投影到当前视角，通过像素级最近邻搜索和深度差值阈值（比例阈值 $\rho{=}0.05$）筛选建立对应对 $(i, m_i)$，再以加权求和方式融合对齐的高斯中心坐标 $\mu_g^t(m_i) = \frac{w_l^t(i) \cdot \mu_l^t(i) + w_g^{t-1}(m_i) \cdot \mu_g^{t-1}(m_i)}{w_l^t(i) + w_g^{t-1}(m_i)}$，并通过轻量 GRU 网络聚合局部与全局的潜在特征。



% METHOD: 2408.13770
\paragraph{TranSplat（深度感知可变形匹配 Transformer）（2024.08）}
TranSplat~\cite{2408_13770} 针对多视角特征匹配不准，尤其非重叠和相似纹理区域。它引入 Depth-aware Deformable Matching Transformer：先用平面扫描得到粗深度分布，再生成深度置信图引导可变形交叉注意力做精匹配。

DDMT 先由 CNN/Transformer 提取多视图特征，并把相机参数和 Depth Anything V2 这类单目深度先验一并编码；粗阶段通过平面扫描计算候选深度分布，细阶段则用深度置信图决定可变形注意力应从哪些局部位置采样。对视图非重叠区域，网络不会强行做不存在的跨视角匹配，而是用单目深度先验补足。最终，每个像素按融合后的深度反投影成高斯中心，并预测尺度、颜色和不透明度。相对 MVSGaussian 的一般 MVS 聚合，TranSplat 把“匹配是否可信”显式接入采样位置和深度融合。
\textbf{设计动机与边界：} 相对原始 3DGS 以独立高斯和图像重建误差为主的设定，TranSplat 把问题落在以下可实现的改动上：TranSplat基于Transformer架构实现可泛化稀疏视角场景重建。核心模块DDMT(Depth-aware Deformable Matching Transformer)：(1)特征提取：CNN+Transformer提取多视角特征，融入相机参数和DepthAnythingV2的单目深度先验；(2)粗匹配：通过平面扫描构建深度候选，计算视角间特征点积获取初始深度分布；(3)粗到精匹配：利用初始深度分布生成深度置信度图，引导可变形交叉注意力进行精确匹配；(4)非重叠区域利用单目深度先验补充。通过置信度加权融合匹配深度和单目深度。每个像素投影为3D Gaussian原语，单次前向传播完成重建。其设计目的不是单独增加一个模块，而是用上述表示、监督或渲染约束消除原流程在目标任务中的具体失配。提出TranSplat，通过深度置信度图引导精确局部特征匹配，利用单目深度估计先验提升非重叠区域深度精度，在RealEstate10K和ACID基准上取得最佳性能。适用时还应保留论文给出的边界：依赖DepthAnythingV2的单目深度先验质量；对极端稀疏(单视角)场景仍有挑战；推理速度受Transformer计算量限制。
\textbf{实现细节：} 与 DepthSplat 在双分支架构中冻结 Depth Anything V2 骨干提取单目特征不同，TranSplat（2024.08） 聚焦于多视角特征匹配本身的精度问题——尤其是在包含大量相似区域和视角间非重叠区域的场景中，纯特征点积匹配容易产生误匹配。TranSplat 的核心思路是两步走：先用深度置信度图引导精确的局部匹配，再用单目深度先验修补非重叠区域的深度。粗到精匹配管线。特征提取阶段采用与 DepthSplat 类似的 CNN 加 Transformer 架构（Swin Transformer 的局部窗口注意力实现跨视角信息交换），并通过 Squeeze-Excitation 层将相机参数注入 CNN 特征提供全局空间信息。同时，从 DepthAnythingV2（base 规格）提取最后一层特征作为深度先验 $D^i \in \mathbb{R}^{H/4 \times W/4 \times C}$ 和相对深度输出。匹配管线分为两个阶段：粗匹配阶段通过平面扫掠立体（$D{=}128$ 个深度候选）和特征点积获取初始深度分布 $d^i_{\mathrm{coarse}}$；精匹配阶段引入深度感知可变形匹配 Transformer（Depth-Aware Deformable Matching Transformer, DDMT），用可变形采样模块（采样 $P{=}4$ 个可变形点）聚合局部空间信息，再由 DMT 根据融合了粗深度、深度特征和相机参数的深度感知特征，自适应地预测采样偏移和深度感知注意力权重 $W^i$，生成深度残差 $\Delta d^i$ 进行精细修正：$d^i_{\mathrm{fine}} = d^i_{\mathrm{coarse}} + \Delta d^i$。直觉上，可变形采样使网络不再仅依赖逐像素对齐的特征点积，而是在局部邻域内搜索更优的匹配位置；DMT 则通过注意力机制将"哪里值得仔细看"的判断融入深度估计过程。Depth Refine U-Net。粗到精匹配仅覆盖视角间有重叠的区域。对于非重叠区域，TranSplat 设计了 Depth Refine U-Net，将精确的多视角匹配几何一致性与 DepthAnythingV2 的可靠相对深度先验融合：以输入图像和相对深度图作为输入，通过 UNet 架构优化深度分布。
\textbf{方法信息：} 代码：\url{https://github.com/xingyoujun/transplat}；引用很多，Semantic Scholar 约 76 次。



% METHOD: 2410.06245
\paragraph{HiSplat（层次化多尺度高斯）（2024.10）}
\textbf{锚点细化：} DepthSplat 和 TranSplat 均生成单一尺度的逐像素高斯，所有高斯具有相同的空间分辨率。HiSplat（2024.10）~\cite{2410_06245}观察到这种单尺度表示存在固有局限：无法同时捕捉大尺度结构（如墙面、地板的整体走向）和小尺度纹理细节（如花纹、文字），导致重建结果中出现模糊、裂纹和定位偏移。其核心思路是引入粗到细的层次化策略，生成多尺度高斯基元并通过跨尺度交互模块协同优化。
HiSplat~\cite{2410_06245} 处理的是单尺度前馈高斯的表达瓶颈。DepthSplat、MVSplat 等方法通常在同一空间分辨率上为像素预测高斯，但大结构需要粗粒度稳定覆盖，纹理和边缘又需要细粒度补偿；统一尺度会让墙面、地面这类大平面定位不稳，也会让小纹理变糊。HiSplat 的改变是把前馈预测做成粗到细的层次化高斯生成。

框架先用 U-Net CNN、多视图 Transformer 和冻结 DINOv2 特征提取三尺度跨视图特征。最低尺度通过代价体预测深度和高斯参数，形成场景粗结构；更高尺度不简单重复预测，而是由 Error Aware Module 读取低尺度的深度误差和图像误差，判断哪些位置需要补偿性高斯；Modulating Fusion Module 再调节多尺度高斯的参数并融合成最终表示。训练时各尺度可以独立渲染并受目标视角监督，推理时使用融合后的层次高斯。

这样，HiSplat 相当于把 ``一次性像素高斯预测'' 改成 ``多分辨率纠错''：低尺度负责稳定几何，高尺度负责修复误差和细节。它的跨尺度融合还避免了每个尺度各自渲染、最后简单相加造成的冲突，而是先调制不同尺度高斯的属性，再输出统一集合。论文在 RealEstate10K、ACID、DL3DV、MipNeRF 360 等数据集上报告相对 pixelSplat/MVSplat 的提升，并强调两视图输入下的稀疏重建能力。\textbf{方法信息：} 代码：\url{https://github.com/open3dvlab/hisplat}；仓库含 demo、检查点和数据处理脚本；GitHub 约 125 stars，Semantic Scholar 约 52 次引用。



% METHOD: 2410.13862
\paragraph{DepthSplat（深度与 Gaussian Splatting 的双向连接）（2024.10）}
\textbf{锚点细化：} DepthSplat：针对 MVSGaussian 纯 MVS 代价体在遮挡、弱纹理、反射面等区域匹配失效的局限，用``MVS 分支 + 冻结 Depth Anything V2 单目分支''的双分支深度模型替代单一 MVS 代价体，核心思想是让多视角匹配提供尺度一致约束、单目先验补充困难区域几何，两个分支互补、缺一不可，同时发现高斯 splatting 渲染损失可作为深度模型的无监督预训练目标，深度估计与视角合成互相增益。
DepthSplat~\cite{2410_13862} 可以作为 ``MVS 深度估计 $\rightarrow$ 高斯中心'' 这条路线的主要锚点。MVSplat 一类方法通常先在多视角之间构造代价体，再把预测深度反投影为像素对齐高斯；问题是遮挡、弱纹理、反射面和大视差区域很容易让纯多视角匹配失败。DepthSplat 的改变不是另起一个高斯表示，而是把预训练单目深度模型的特征接入多视图匹配分支，让单目先验补困难区域，多视角匹配补尺度一致性。

具体流程是：输入带位姿的多视角图像后，模型一边提取多视角匹配特征并构造代价体，一边引入单目深度特征；融合后的深度分支输出每个输入像素的深度，再反投影成 3D 高斯中心。随后一个轻量 Gaussian adapter 预测尺度、旋转、不透明度和颜色，目标视角仍通过 3DGS rasterization 渲染。因为高斯中心、深度图和目标视角渲染共享同一几何链条，深度分支的误差不会被隐藏在颜色头里，而会在重投影图像中体现。DepthSplat 还有反向使用方式：把 splatting 渲染误差作为无监督预训练目标来训练深度模型，因此它同时回答了两个问题，即深度先验如何提升前馈 3DGS，以及 3DGS 渲染如何反过来提供深度监督。

这篇工作的重要性在于把 ``深度 foundation model'' 引入前馈 3DGS 的几何主干，而不是只把深度当作后处理约束。它把高斯中心预测转化为密集深度估计问题，并让渲染损失直接暴露深度错误：如果深度偏移，反投影高斯会在目标视角产生可见错位。论文报告 12 张 $512\times960$ 输入视图约 0.6 秒完成前馈重建，并在 ScanNet、RealEstate10K、DL3DV 上验证深度估计和新视角合成的互补收益。\textbf{方法信息：} Microsoft 参与了这篇论文；代码：\url{https://github.com/cvg/depthsplat}；仓库提供模型、数据说明、渲染脚本和保存高斯为 PLY 的选项；GitHub 约 1216 stars，Semantic Scholar 约 260 次引用。



% METHOD: 2411.04924
\paragraph{MVSplat360（潜空间特征与视频扩散 360 合成）（2024.11）}
\textbf{锚点细化：} ProSplat 利用一步扩散模型增强单个渲染视角，但无法保证多个增强视角之间的时序一致性——这对于 360\textdegree{} 全景合成尤为关键。MVSplat360（2024.11）~\cite{2411_04924}选择预训练的 Stable Video Diffusion（SVD）模型作为增强骨干——SVD 在视频数据上训练，天然具备时序一致性约束。
MVSplat360~\cite{2411_04924} 把前馈 3DGS 推到宽范围甚至 360 度新视角合成。这个任务比常规插值更难，因为 5 张左右输入图像之间可能几乎没有重叠，纯几何前馈模型即使能估计粗结构，也会在未观测区域缺外观。MVSplat360 的改变是让 3DGS 不直接渲染 RGB，而是渲染可供视频扩散模型使用的潜空间特征。

第一阶段仍以 MVSplat 式粗几何为基础：跨视角 Transformer 融合稀疏输入，代价体估计深度，并预测 3D 高斯参数。不同之处在于，每个高斯还携带 latent feature，目标视角 rasterization 输出的不是颜色图，而是 Stable Video Diffusion 可读的潜特征图。第二阶段，SVD 以这些特征作为姿态和视觉条件执行去噪，利用视频生成模型的时序一致性生成一组连续目标视角。几何分支给扩散模型提供相机约束和粗外观线索，扩散分支补全前馈几何无法覆盖的细节。

因此，相对普通 MVSplat 直接渲染 RGB 的做法，MVSplat360 不是简单的渲染后增强，而是把 ``前馈高斯粗几何'' 与 ``多帧一致外观生成'' 接成端到端系统。潜特征仍由 3DGS 按目标相机投影得到，所以扩散模型看到的条件图携带几何可见性和视角变化，而不是无约束的文本或单帧图像提示。它在 DL3DV-10K 建立 360 度 NVS 基准，并在 RealEstate10K 等数据上比较宽范围视角合成。\textbf{方法信息：} 代码：\url{https://github.com/donydchen/mvsplat360}；仓库提供 SVD 权重接入、DL3DV/RealEstate10K 数据说明和评测脚本；GitHub 约 315 stars，Semantic Scholar 约 109 次引用。
\textbf{设计动机与边界：} 相对原始 3DGS 以独立高斯和图像重建误差为主的设定，MVSplat360 把问题落在以下可实现的改动上：框架包含两个模块：(1) 多视角粗几何重建——基于MVSplat前馈3DGS模型，从稀疏视图预测3D高斯参数。关键改进是在预测高斯属性的同时预测额外的高斯特征，可光栅化为潜空间特征而非RGB图像。使用交叉视图Transformer编码器融合多视图信息，构建代价体积进行深度估计，改进视角选择策略增强对大间距输入的鲁棒性；(2) 多帧一致外观精化——将粗几何重建模块渲染的潜空间特征作为Stable Video Diffusion（SVD）的条件输入，引导去噪过程生成照片级、3D一致的新视角。整个框架端到端可训练，无需逐场景优化。其设计目的不是单独增加一个模块，而是用上述表示、监督或渲染约束消除原流程在目标任务中的具体失配。提出MVSplat360——首个前馈360度场景合成方法。核心创新：(1) 将前馈3DGS模型（基于MVSplat）重构为在SVD潜空间渲染特征而非RGB图像；(2) 渲染的潜空间特征作为姿态和视觉线索引导SVD去噪过程，生成照片级3D一致视图；(3) 端到端可训练，支持仅5张稀疏输入视图渲染任意视角。在DL3DV-10K上建立新基准。适用时还应保留论文给出的边界：极度稀疏视图下几何重建仍可能不完整；依赖预训练SVD模型的质量；处理时间受扩散模型影响。
% METHOD: 2411.06390
\paragraph{SplatFormer（点 Transformer 精炼高斯集合）（2024.11）}
SplatFormer~\cite{2411_06390} 不从图像直接预测高斯，而是接收已有 3DGS 集合并用 Point Transformer 做一次前向精炼。它关注训练视角分布外的新视角伪影：有限视角优化出的高斯在极端新视角下会暴露结构错误。

模型把高斯参数视为点集合，输入包括位置、尺度、旋转、不透明度和颜色等属性；Point Transformer 的局部/全局注意力建模基元之间的空间关系，并在单次前向中输出修正后的高斯属性。它不要求固定数量的输入视角，也不通过扩散模型重新生成图像，而是直接修正可渲染资产本身。相对正则化或后处理，SplatFormer 的锚点是“在高斯集合上学习一个泛化修复器”，因此特别适合极端新视角下的 OOD 伪影处理。
\textbf{实现细节：} 前述方法均在二维图像空间或潜在空间操作扩散模型，SplatFormer（2024.11） 选择了一条根本不同的路线：直接在三维高斯基元集合上操作。这是首个为高斯 splat 设计的点 Transformer 模型。SplatFormer 接收有限训练视角优化得到的初始 3DGS 集合作为输入（而非二维图像），通过单次前向传播直接精炼高斯参数——每个高斯被视为一个"点"，其属性（位置、尺度、旋转、颜色、不透明度）作为点特征。Transformer 的自注意力在三维高斯之间建立全局交互：精炼过程中，一个高斯的属性调整可以参考空间中其他所有高斯的状态，实现全局协调的修复——例如，消除在分布外视角下暴露的"针刺"伪影需要同时调整多个空间上相邻的高斯的尺度。
\textbf{设计动机与边界：} 相对原始 3DGS 以独立高斯和图像重建误差为主的设定，SplatFormer 把问题落在以下可实现的改动上：SplatFormer的核心创新在于将点Transformer架构直接应用于3DGS集合。传统方法通过正则化或数据驱动先验来改善泛化，但效果有限。SplatFormer接收有限训练视角下优化的初始3DGS集合作为输入，利用点Transformer的自注意力机制建模高斯基元之间的关系，在单次前向传播中精炼高斯属性以消除OOD视角的伪影。这是首个成功将点Transformer直接应用于3DGS集合的工作，突破了以往多场景训练方法在推理时只能处理有限输入视角数量的限制。模型在多种合成和真实数据集上进行了全面评估。其设计目的不是单独增加一个模块，而是用上述表示、监督或渲染约束消除原流程在目标任务中的具体失配。首次将点Transformer直接应用于3DGS集合的精炼，系统性地评估了3DGS及相关方法在OOD测试视角下的表现，提出的方法在极端新颖视角下实现了SOTA渲染质量，显著超越各种3DGS正则化技术和扩散模型框架。适用时还应保留论文给出的边界：依赖初始3DGS的质量；推理时需要额外的Transformer前向传播开销；模型在未见过的场景类型上可能泛化有限
\textbf{方法信息：} 代码：\url{https://github.com/chenyutongthu/splatformer}；GitHub 关注度较高，约 405 stars。



% METHOD: 2503.20168
\paragraph{EVolSplat（城市场景体素前馈）（2025.03）}
EVolSplat~\cite{2503_20168} 面向城市街景。它先用 Metric3D/UniDepth 等深度模型生成全局点云，并用深度一致性和 3D 滤波去除噪声，再把点云转成统一 3D 体素。

3D CNN 在体素中预测几何属性，IBR 模块从聚合 2D 特征预测高频颜色；远处天空和背景用半球背景高斯建模。相对 VolSplat 的通用体素框架，EVolSplat 更针对城市前景/远景分解。
全局点云初始化是它区别于普通像素高斯的第一步：Metric3D/UniDepth 给出度量深度，深度一致性检查和 3D 滤波先清理不一致点，再反投影到世界坐标。点云体素化后，3D CNN 在局部体积上下文中预测位置偏移、尺度、旋转和不透明度，把带噪点云变成可渲染高斯。外观方面，3D CNN 不擅长恢复图像高频纹理，因此 EVolSplat 另接遮挡感知 IBR 模块，从聚合 2D 特征预测颜色；远景和天空则由半球背景模型承担。
\textbf{方法信息：} Huawei 参与了这篇论文。
\textbf{设计动机与边界：} 相对原始 3DGS 以独立高斯和图像重建误差为主的设定，EVolSplat 把问题落在以下可实现的改动上：EVolSplat的核心设计包括三个组件。（1）全局点云初始化：利用预训练深度模型（Metric3D/Unidepth）预测度量深度，通过深度一致性检查和3D滤波去除不一致值和浮点伪影，将单帧深度反投影到全局世界坐标构建场景点云。（2）3D CNN编码体素：将全局点云转化为3D体素，通过3D卷积网络解码几何属性（位置偏移、尺度、旋转、不透明度），将噪声点云转化为3D高斯。3D CNN能从3D上下文学习准确的几何预测。（3）IBR模块和半球背景模型：3D CNN难以捕捉高频外观细节，因此引入遮挡感知的基于图像的渲染模块从聚合2D特征预测高斯颜色。对于远处视图和天空，使用位于远处半球上的3D高斯进行建模。整个场景被分解为前景（体素范围内）和背景（远处景色和天空）两部分独立建模。其设计目的不是单独增加一个模块，而是用上述表示、监督或渲染约束消除原流程在目标任务中的具体失配。提出了EVolSplat，一种面向城市场景的高效前馈3DGS模型，使用3D CNN在统一体素中预测多帧3D高斯，解决了现有像素对齐方法的多视角不一致和内容重复问题，实现了快速前馈重建和实时渲染。适用时还应保留论文给出的边界：3D CNN的平滑偏差影响高频细节；前景体素范围需要预定义；深度预测质量影响初始化效果。
% METHOD: 2503.22986
\paragraph{FreeSplat++（大规模室内场景增量融合）（2025.03）}
FreeSplat++~\cite{2503_22986} 将 FreeSplat 扩展到大规模室内全场景重建。它保留低成本跨视角聚合和 PTF，但把融合做成增量式全局过程，类似把传统 TSDF 融合思想迁移到高斯三元组。

局部预测得到的高斯三元组先由纯 CNN 骨干和自适应 3-point 代价体估计深度，再与已有全局三元组对齐；重叠区域不简单拼接，而是在潜特征空间融合冗余高斯并累积跨视角证据。Weighted Floater Removal 利用融合权重降低漂浮高斯不透明度，因为重复视图都支持的基元更可能是真实表面，只被少数视图支持的基元更可能是浮点。最后，系统用前馈阶段得到的一致深度图做逐场景微调，在不破坏几何的情况下提升渲染质量。相对 FreeSplat，它的增量是把“少量输入的前馈预测”推进到“长序列室内场景的可增量融合系统”。
\textbf{设计动机与边界：} 相对原始 3DGS 以独立高斯和图像重建误差为主的设定，FreeSplat++ 把问题落在以下可实现的改动上：FreeSplat++包含四个核心组件：(1)低成本跨视角聚合：使用纯CNN骨干网络处理长序列高分辨率输入，自适应构建3pt代价体，通过平面扫描立体匹配获取深度；(2)像素对齐三元组融合（PTF）：受传统TSDF融合启发，将全局高斯三元组与局部三元组增量对齐，在潜在特征空间融合冗余3D高斯并聚合跨视角特征；(3)加权浮点去除（WFR）：利用PTF过程中累积的权重调整浮点高斯的不透明度；(4)深度正则化逐场景微调：利用前馈预测阶段获得的多视角一致深度图作为额外约束，微调整个场景的3DGS以提升渲染质量同时保持几何精度。其设计目的不是单独增加一个模块，而是用上述表示、监督或渲染约束消除原流程在目标任务中的具体失配。将可泛化3DGS扩展为大规模室内全场景重建的替代方案，提出PTF模块增量聚合多视角重叠高斯基元以自适应减少冗余，WFR策略有效去除浮点，深度正则化微调在保持几何精度的同时提升渲染质量。适用时还应保留论文给出的边界：依赖CNN骨干网络的表达能力；对极端视角变化和纹理缺失区域的处理仍有限；深度正则化微调需要额外时间
\textbf{实现细节：} FreeSplat++（2025.03） 将 FreeSplat 的室内大场景重建扩展到全场景级别，解决三个核心问题：长序列处理效率、跨视角浮游体（floater，悬浮在场景空间中的半透明碎片状伪影）噪声、以及前馈质量与逐场景优化之间的差距。与 FreeSplat 的关键差异。(1) FreeSplat++ 将解投影分辨率降至全分辨率的 $1/2$（而非 FreeSplat 的全分辨率），使全场景重建中的高斯数量保持在可承受范围内，同时性能几乎不受影响。(2) 提出加权浮游体去除（Weighted Floater Removal, WFR）模块。该模块利用 PTF 过程中累积的融合权重作为可靠度指标：权重过低的高斯表示其在多视角融合中被很少观测到，极可能是浮游体噪声。WFR 据此降低这些高斯的不透明度，作为显式深度融合方法有效抑制噪声深度图引入的浮游体。(3) 设计深度正则化的逐场景微调：利用前馈预测阶段获得的多视角一致深度图作为额外约束，在保持几何精度的同时精炼高斯参数的渲染质量。在 ScanNet 全场景实验中，FreeSplat++ 的前馈预测已在几何精度上优于传统逐场景优化 3DGS；经 10K 次迭代微调后，渲染质量显著提升，而训练时间仅为传统 3DGS（30K 次迭代）的 $1/3$。
\textbf{方法信息：} 代码：\url{https://github.com/wangys16/freesplat}；GitHub 关注度较高，约 164 stars。



% METHOD: 2504.09588
\paragraph{TextSplat（文本引导语义融合）（2025.04）}
TextSplat~\cite{2504_09588} 在可泛化 3DGS 中加入文本条件。它并行使用扩散先验深度估计器、语义分割网络和多视角交互网络，再用文本引导的注意力融合三类表征。

CLIP 文本嵌入作为跨模态条件，帮助模型聚合对象级语义、全局布局和多视角几何。相对 Uni3R/GSemSplat 的语义重建，TextSplat 的特点是文本提示直接参与高斯参数预测过程。
DPDE 分支用文本条件扩散先验形成深度表征，SASN 分支提取对象级语义和全局场景布局，多视角交互网络则用 Transformer 处理跨视角匹配。融合阶段不是简单拼接三类特征，而是让多级 CLIP 文本嵌入通过注意力选择哪些深度、语义和几何线索应进入高斯参数预测。这样，文本不只是推理时查询语义场，而是在重建时就影响哪些结构和对象关系被强调。
\textbf{实现细节：} TextSplat（2025.04） 是首个文本驱动的可泛化高斯 splatting 框架，其核心思想是利用文本描述作为语义桥梁，引导多模态信息的融合以增强三维重建质量。框架包含三个并行模块获取互补表示：扩散先验深度估计器（利用预训练扩散模型的几何先验提供准确的深度信息）、语义感知分割网络（从输入图像提取细粒度语义信息）、多视角交互网络（通过跨视角特征交互精炼多视角表示）。三种表示在文本引导语义融合模块（Text-Guided Semantic Fusion Module）中通过注意力机制集成——文本嵌入引导网络关注与语义描述相关的图像区域，使融合后的三维高斯参数同时编码丰富的几何、外观和语义线索。实验表明，文本引导的语义融合显著提升了复杂场景中细粒度细节的重建质量。
\textbf{方法信息：} ByteDance 参与了这篇论文。
\textbf{设计动机与边界：} 相对原始 3DGS 以独立高斯和图像重建误差为主的设定，TextSplat 把问题落在以下可实现的改动上：TextSplat框架包含三个并行模块和一个融合模块：(1) 扩散先验深度估计器（DPDE）：利用文本条件生成先验桥接2D观测与3D理解，通过去噪UNet获取扩散深度表征。(2) 语义感知分割网络（SASN）：编码器-解码器架构捕获多尺度语义线索，提供对象级语义和全局场景布局。(3) 多视角交互网络：利用跨视角Transformer增强特征匹配。(4) 文本引导语义融合模块：通过多级文本嵌入的注意力机制整合上述三种表征，生成增强的3D高斯参数。文本提示通过CLIP编码为跨模态特征，引导特征聚合过程。其设计目的不是单独增加一个模块，而是用上述表示、监督或渲染约束消除原流程在目标任务中的具体失配。首个文本驱动的可泛化高斯泼溅框架，通过融合扩散先验深度、语义分割和多视角交互三种互补表征，利用文本引导的注意力机制聚合特征，提升3D重建质量。适用时还应保留论文给出的边界：依赖预训练文本编码器的质量；对复杂场景描述的处理能力有限；增加的计算开销来自多个并行模块。
% METHOD: 2505.15185
\paragraph{MonoSplat（单目深度基础模型适配）（2025.05）}
MonoSplat~\cite{2505_15185} 用冻结单目深度基础模型提供强视觉先验。它从 Depth Anything 编码器提取多尺度特征，再用 Mono-Multi adapter 将单目特征转换为跨视角一致的多视图表示。

适配器通过多视图 Transformer 交换特征，集成高斯预测模块融合单目深度先验和多视图匹配特征，预测位置、不透明度、协方差和颜色。相对 DepthSplat，它更强调冻结 depth foundation model，只训练轻量适配层以保持泛化。
\textbf{实现细节：} 与 DepthSplat 直接冻结预训练单目深度骨干（Depth Anything V2 的 ViT）提取单目特征的策略不同，MonoSplat（2025.05） 提出了一种更灵活的单目先验利用方式：不冻结骨干，而是设计轻量级的单目-多视角特征适配器（Mono-Multi Feature Adapter），以最少的可训练参数将单目特征转换为多视角表示。这一设计的核心优势在于效率：适配器的可训练参数远小于从头训练整个骨干，却在跨场景泛化上表现卓越。其动机来自一个实际观察：预训练单目深度模型在大规模数据上学到的视觉先验极为丰富（能处理各种光照、纹理、遮挡条件），但直接冻结使用时这些先验被"锁死"在单目预测范式中——模型只见过从单张图像预测深度的任务，不天然具备多视角几何推理能力。MonoSplat 的适配器通过轻量注意力机制实现两个目标：(1) 跨视角对齐——将不同视角的单目特征映射到统一的多视角特征空间中；(2) 先验保持——在适配过程中保留有价值的单目几何先验不被破坏。适配后的多视角特征与原始单目特征在集成高斯预测模块（Integrated Gaussian Prediction）中融合，共同解码高斯基元的几何和外观参数。
\textbf{方法信息：} 代码：\url{https://github.com/cuhk-aim-group/monosplat}。
\textbf{设计动机与边界：} 相对原始 3DGS 以独立高斯和图像重建误差为主的设定，MonoSplat 把问题落在以下可实现的改动上：MonoSplat框架包含两个核心组件：（1）Mono-Multi特征适配器：第一阶段从冻结的Depth Anything编码器中提取多尺度特征并统一聚合；第二阶段利用多视图Transformer进行跨视图特征交换，建立几何一致性。同时提取互补的单目特征提供丰富的深度先验。（2）集成高斯预测模块：协同结合单目和多视图特征预测精确的高斯基元（位置、不透明度、协方差、球谐颜色）。冻结深度编码器保持泛化能力，仅训练轻量级适配层和预测模块，可训练参数极少。整个框架通过前馈网络直接从图像预测3D高斯，无需逐场景优化。其设计目的不是单独增加一个模块，而是用上述表示、监督或渲染约束消除原流程在目标任务中的具体失配。提出利用预训练单目深度基础模型的丰富视觉先验进行鲁棒高斯重建。设计了Mono-Multi特征适配器将单目特征转换为多视图表示，以及集成高斯预测模块融合两类特征进行精确高斯基元生成。冻结深度编码器，仅训练轻量级适配器和预测模块，在保持计算效率的同时实现卓越的泛化能力。适用时还应保留论文给出的边界：深度编码器冻结可能限制在某些特定域的适应性；多视图数量较少时性能可能下降；对极端光照变化的鲁棒性有待验证。
% METHOD: 2506.09479
\paragraph{TinySplat（前馈高斯紧凑表示）（2025.06）}
TinySplat~\cite{2506_09479} 接在前馈 3DGS 输出之后，解决前馈模型生成高斯过多的问题。它不重新训练预测网络，而是提供一套无训练压缩流程。

View-Projection Transformation 利用前馈高斯和像素一一对应的结构，把几何参数投影到更紧凑空间；Visibility-Aware Basis Reduction 将外观特征能量沿主导可见方向集中，减少感知冗余；剩余空间冗余交给视频 codec。相对 EcoSplat 的条件化模型内控制，TinySplat 是后处理压缩。
\textbf{实现细节：} EcoSplat 聚焦于基元数量控制，而前馈 3DGS 输出的存储压缩同样至关重要。TinySplat（2025.06） 提出完整的前馈压缩框架，关注的是存储压缩——如何在保持渲染质量的前提下将前馈 3DGS 的数据体积压缩到极致。其核心洞察是前馈 3DGS 生成的像素对齐高斯特征图中蕴含三类不同性质的冗余，可分别消除：(1) 几何冗余——通过视角投影变换（View-Projection Transformation, VPT）消除：VPT 将高斯几何参数从世界坐标系变换到相机空间，利用像素对齐高斯的空间正则性仅编码相对于像素中心的偏移量，并近似地将旋转和缩放参数变换至相机坐标系以增强局部相关性，使几何参数在新的坐标空间中呈现更紧凑的分布。(2) 感知冗余——通过可见性感知基约简（Visibility-Aware Basis Reduction, VABR）消除：VABR 根据各球谐基函数在有效观测方向上的平均响应值计算可见性权重，通过加权协方差矩阵的特征分解提取主成分基，将高维 SH 系数投影到紧凑的可见性感知子空间中——仅保留对渲染贡献最大的方向分量。(3) 空间冗余——通过现成视频编解码器处理。整个压缩框架无需训练，直接应用于前馈 3DGS 的输出。
\textbf{方法信息：} 代码：\url{https://github.com/maxgillett/tinysplat}。
\textbf{设计动机与边界：} 相对原始 3DGS 以独立高斯和图像重建误差为主的设定，TinySplat 把问题落在以下可实现的改动上：TinySplat基于标准前馈3DGS方法（如MVSplat、DepthSplat）生成高斯原语后，集成无需训练的压缩框架系统消除三类冗余：(1)View-Projection Transformation(VPT)：将几何参数投影到更紧凑的空间以减少几何冗余，利用前馈方法中高斯与像素一一对应的特性；(2)Visibility-Aware Basis Reduction(VABR)：通过基变换将特征能量沿主导视角方向对齐，消除感知冗余，使外观特征在视觉上更关键的方向上集中能量；(3)空间冗余通过现成视频编解码器处理。整个压缩框架无需训练，编码时间仅为SOTA方法的25%，解码时间仅为1%。其设计目的不是单独增加一个模块，而是用上述表示、监督或渲染约束消除原流程在目标任务中的具体失配。提出首个完全前馈的3D场景紧凑表示生成方法，通过View-Projection Transformation(VPT)消除几何冗余、Visibility-Aware Basis Reduction(VABR)消除感知冗余、视频编解码器消除空间冗余，实现100倍以上压缩。适用时还应保留论文给出的边界：压缩质量依赖于前馈方法的初始高斯质量；对非前馈方法生成的3DGS适用性有限；VPT和VABR参数需根据场景调整。
% METHOD: 2507.14921
\paragraph{Stereo-GS（几何与外观解耦的立体前馈）（2025.07）}
Stereo-GS~\cite{2507_14921} 把 3DGS 预测拆成几何点图和外观高斯特征两个任务。它先用预训练扩散模型从单张输入生成多视角图像，再用 Siamese ViT 和跨视角注意力提取立体 token。

所有立体对的 token 被拼成多视角立体 token，并由全局自注意力预测统一世界坐标系下的 point map。几何头负责点图，外观头负责高斯颜色和特征，二者组合为 GS-map 后再由细化网络提高渲染质量。这个解耦很关键：几何可以借助预训练立体模型形成稳定坐标，外观再围绕该几何补充视角相关颜色。相对 DepthSplat 的深度回归，它更接近多视角立体基础模型，把位姿无关重建拆成“先稳定几何、再补高斯外观”。
\textbf{设计动机与边界：} 相对原始 3DGS 以独立高斯和图像重建误差为主的设定，Stereo-GS 把问题落在以下可实现的改动上：Stereo-GS的核心流程：(1)多视图图像生成：利用预训练扩散模型从单张输入图像生成多视角图像；(2)局部立体特征提取：使用Siamese ViT编码器提取图像对token，通过双胞胎ViT解码器进行交叉注意力获取立体特征token；(3)多视角全局融合：将所有立体对的token拼接为多视角立体token，通过多视角自注意力块学习全局表示并预测统一世界坐标系下的3D点图；(4)解耦预测：专用点预测头生成多视角点图（几何），高斯预测头生成高斯特征（外观），组合为GS-map；(5)细化网络：增强GS-map以实现高质量重建。关键创新是将立体视觉模型从双视角提升到多视角，实现前馈式统一坐标系的3D重建。其设计目的不是单独增加一个模块，而是用上述表示、监督或渲染约束消除原流程在目标任务中的具体失配。提出Stereo-GS，首次将3DGS预测解耦为几何（点图）和外观（高斯特征）两个独立任务。利用预训练立体视觉模型作为几何先验减少对数据驱动先验的依赖，实现无相机参数的位姿无关3D重建，显著降低了训练资源需求。适用时还应保留论文给出的边界：依赖预训练扩散模型生成多视角图像，引入额外误差；解耦策略可能在复杂几何场景中损失信息；全局注意力的计算复杂度随视角数增加。
\textbf{方法信息：} NVIDIA 参与了这篇论文；代码：\url{https://github.com/sadrasafa/StereoGS}。



% METHOD: 2508.20754
\paragraph{C$^3$-GS（上下文、跨维度、跨尺度特征）（2025.08）}
C$^3$-GS~\cite{2508_20754} 基于 MVSGaussian 加入三个轻量模块。CGA 将坐标感知上下文嵌入 2D 特征，增强长程通道依赖；CDA 融合 2D 特征和 3D 代价体，形成空间感知描述符；CSF 跨尺度更新高斯不透明度。

流程仍是 FPN 特征、代价体、3D CNN 正则化、深度回归和高斯解码。相对 MVSGaussian，它不是换范式，而是补强特征融合，使稀疏输入下的上下文更稳。
CGA 在 2D 特征中加入坐标感知上下文，用注意力捕获长程通道依赖，降低代价匹配对局部纹理的依赖；CDA 把 2D 特征与 3D 代价体融合成空间感知描述符，让跨视角约束不仅停留在图像平面；CSF 在多尺度上更新高斯不透明度，使全局结构和局部细节都能进入最终渲染。整体上，C$^3$-GS 是 MVSGaussian 的特征融合增强版，而不是新的高斯参数化。
\textbf{方法信息：} 代码：\url{https://github.com/yuhsihu/c3-gs}。
\textbf{设计动机与边界：} 相对原始 3DGS 以独立高斯和图像重建误差为主的设定，${C}^{3}$-GS 把问题落在以下可实现的改动上：C³-GS基于MVSGaussian构建，在统一渲染管线中集成三个轻量模块：(1) CGA（上下文感知引导注意力）：在2D特征中嵌入坐标感知上下文，通过两个注意力图捕获长程特征通道依赖关系，增强代价匹配的鲁棒性；(2) CDA（跨维度注意力）：通过融合2D特征与3D代价体构造3D空间感知描述符，学习跨视角和上下文感知约束；(3) CSF（跨尺度融合）：跨尺度自适应更新高斯不透明度，在保持全局结构的同时保留细节。整体流程：FPN提取多尺度特征→CGA增强→可微单应变换构建代价体→3D CNN正则化→深度回归→高斯参数解码→CSF融合渲染。其设计目的不是单独增加一个模块，而是用上述表示、监督或渲染约束消除原流程在目标任务中的具体失配。提出C³-GS框架，通过三个轻量模块增强特征学习：(1) 上下文感知引导注意力（CGA）将坐标感知上下文嵌入2D特征；(2) 跨维度注意力（CDA）融合2D特征与3D体积构建空间感知描述符；(3) 跨尺度融合（CSF）跨尺度精炼高斯不透明度。在四个基准数据集上达到SOTA渲染质量和泛化能力。适用时还应保留论文给出的边界：基于MVSGaussian构建，受限于MVS范式的固有局限；对极端稀疏视角（如仅2视角）的提升空间有限
% METHOD: 2509.19297
\paragraph{VolSplat（体素对齐高斯预测）（2025.09）}
VolSplat~\cite{2509_19297} 反思像素对齐范式。每个输入像素预测一个高斯会造成视图依赖密度、重复基元和对齐误差；VolSplat 的改变是先把多视角特征聚合到统一 3D 体素空间，再从体素特征预测高斯。

方法用 ResNet 提取多视图特征，交叉注意力细化后构建逐视图代价体，再通过平面扫描把候选深度上的特征放进统一 voxel grid。此后，3D U-Net 在体素空间内做几何和外观推理，输出的高斯由体素结构决定，而不是刚性绑定到某个输入像素。这样，多视角信息先在 3D 中融合，再决定哪里需要高斯、需要多少高斯。相对 DepthSplat 的像素深度反投影，VolSplat 更适合处理视图重叠复杂、像素对应不稳定的场景。
\textbf{设计动机与边界：} 相对原始 3DGS 以独立高斯和图像重建误差为主的设定，VolSplat 把问题落在以下可实现的改动上：VolSplat提出将前馈3DGS的预测范式从像素对齐转变为体素对齐：（1）3D特征构建：使用ResNet骨干网络提取多视图特征，通过交叉注意力细化，构建逐视图代价体，然后使用平面扫描策略聚合到统一的3D体素空间，有效将特征融合与相机视锥解耦。（2）3D U-Net推理：在统一体素空间内使用3D U-Net对场景几何和外观进行体积推理，而非在2D特征空间中融合。（3）体素对齐高斯预测：直接从细化的体素特征预测高斯原语，使3D高斯的分布由体积结构本身决定，而非刚性地与2D像素网格绑定。这种方法的优势包括：减少浮点和视图依赖不一致性（因为多视图信息在3D空间中聚合），允许高斯数量根据3D场景复杂度自适应控制，避免像素对齐中离散像素位置采样的误差。其设计目的不是单独增加一个模块，而是用上述表示、监督或渲染约束消除原流程在目标任务中的具体失配。识别了现有前馈3DGS方法中像素对齐预测范式的固有局限（视图依赖密度分布、对齐误差），提出体素对齐高斯预测新范式VolSplat。直接从3D体素网格预测高斯，克服了像素对齐对误差敏感的2D特征匹配的依赖，实现更鲁棒的多视角一致性和自适应密度控制。适用时还应保留论文给出的边界：3D体素网格的分辨率受限于内存；对大尺度场景需要多尺度策略；推理速度可能受3D U-Net影响。
\textbf{实现细节：} 前述方法从 token 和锚点角度脱离像素对齐。VolSplat（2025.09） 用体素对齐取代像素对齐——从预测的三维体素网格直接解码高斯参数，避免了易错的二维特征匹配，实现了基于三维场景复杂度的自适应密度控制。与 GlobalSplat 在潜在空间中用 token 表示全局场景不同，VolSplat 的体素表示保持了显式的三维空间结构。其管线分为三个阶段：(1) 从多视角图像提取二维特征并构建逐视角代价体，通过深度预测模块估计各视角深度图；(2) 利用深度图将二维特征解投影至三维空间，构建体素特征网格，再由稀疏 3D U-Net 在体素空间中精炼特征；(3) 从精炼后的体素特征直接预测每个占据体素的高斯参数。这一设计解决了像素对齐的两个核心局限：其一，像素级采样在遮挡、弱纹理和运动视差区域产生不一致的采样模式，体素聚合则将多视角信息融合到统一的 3D 容器中，消除了视角偏差的密度分布；其二，像素对齐强制均匀密度，体素对齐则允许根据体素内三维场景的复杂度自适应分配高斯密度——场景结构复杂的体素获得更多基元，平坦区域则更稀疏。在 RealEstate10K、ScanNet 和 ACID 上的实验证实了 VolSplat 在几何一致性和渲染质量上的一致优势。
\textbf{技术实现补充：} 前述方法从 token 和锚点角度脱离像素对齐。VolSplat（2025.09） 用体素对齐取代像素对齐——从预测的三维体素网格直接解码高斯参数，避免了易错的二维特征匹配，实现了基于三维场景复杂度的自适应密度控制。与 GlobalSplat 在潜在空间中用 token 表示全局场景不同，VolSplat 的体素表示保持了显式的三维空间结构。其管线分为三个阶段：(1) 从多视角图像提取二维特征并构建逐视角代价体，通过深度预测模块估计各视角深度图；(2) 利用深度图将二维特征解投影至三维空间，构建体素特征网格，再由稀疏 3D U-Net 在体素空间中精炼特征；(3) 从精炼后的体素特征直接预测每个占据体素的高斯参数。这一设计解决了像素对齐的两个核心局限：其一，像素级采样在遮挡、弱纹理和运动视差区域产生不一致的采样模式，体素聚合则将多视角信息融合到统一的 3D 容器中，消除了视角偏差的密度分布；其二，像素对齐强制均匀密度，体素对齐则允许根据体素内三维场景的复杂度自适应分配高斯密度——场景结构复杂的体素获得更多基元，平坦区域则更稀疏。在 RealEstate10K、ScanNet 和 ACID 上的实验证实了 VolSplat 在几何一致性和渲染质量上的一致优势。
\textbf{方法信息：} 代码：\url{https://github.com/ziplab/VolSplat}；GitHub 关注度较高，约 288 stars。



% METHOD: 2510.08575
\paragraph{ReSplat（递归误差反馈前馈重建）（2025.10）}
\textbf{锚点细化：} ReSplat：针对 DepthSplat 单步前馈受网络容量限制无法反复改进的局限，用``渲染误差驱动的循环迭代精炼''替代一次性推断，核心思想是在 $16\times$ 下采样空间预测紧凑初始高斯（像素对齐方法的 $1/16$），每轮用当前高斯渲染输入视角并与真实图像比较，像素级和特征级误差经全局注意力聚合后驱动 $k$NN 注意力块预测逐高斯残差更新——所谓 $k$NN 注意力即每个高斯仅与三维空间中 $k$ 个最近邻高斯交互而非全部高斯，符合高斯更新本质上是局部操作（移动位置、调整大小）的特性，整个过程为纯前向推断、不涉及梯度回传，以 $1/16$ 高斯数量超越像素对齐方法。
ReSplat~\cite{2510_08575} 直接继承 DepthSplat 的深度-高斯前馈框架，但认为单次前向推断仍然太硬：网络一次输出高斯后便不再查看自己的渲染错误，测试场景分布、输入视角数或分辨率稍有变化时，很难像逐场景优化那样继续修正。它的增量是把 ``渲染误差'' 变成一个可学习的反馈信号，用循环网络做无梯度精化。

初始阶段，ReSplat 不再为每个输入像素预测高斯，而是在 $16\times$ 下采样空间生成紧凑高斯集合，再用交替的 $k$NN 注意力和全局注意力编码 3D 上下文。递归阶段，每一轮先把当前高斯渲染回输入视角，得到渲染图、真实图与误差图；这些图像特征被送入共享权重的更新网络，输出高斯中心、尺度、不透明度等参数的残差更新。这里的 $k$NN 注意力只让一个高斯与空间近邻交互，适合高斯局部移动和尺度调整；全局注意力则补充远距离结构关系。整个测试过程是多次前向传播，不需要对测试场景做反向传播。

因此，ReSplat 在技术位置上不是替代 DepthSplat 的几何初始化，而是给前馈模型加上一个学到的优化器：先用少量高斯快速给出场景骨架，再让渲染误差逐轮指出哪里缺细节、哪里颜色或几何不一致。由于更新网络权重共享，递归步数可以在质量和耗时之间调节；误差图来自当前测试场景本身，所以反馈信号比固定训练先验更贴近输入分布。论文在 DL3DV、RealEstate10K 和 ACID 上报告跨 2、8、16、32 输入视角和多种分辨率的泛化结果。\textbf{方法信息：} Microsoft 参与了这篇论文；代码：\url{https://github.com/cvg/resplat}；仓库 README 明确提供预训练模型、COLMAP demo、评测脚本和两阶段训练脚本；GitHub 约 284 stars。
\textbf{设计动机与边界：} 相对原始 3DGS 以独立高斯和图像重建误差为主的设定，ReSplat 把问题落在以下可实现的改动上：ReSplat包含两个核心阶段：（1）初始高斯重建：基于DepthSplat架构，但预测16倍下采样空间的高斯（先4倍空间压缩预测深度图，再unproject到3D获得M=N×16个点）。为补偿信息损失，使用6个交替的kNN注意力和全局注意力层编码3D上下文。（2）递归高斯更新：训练权重共享的递归网络，利用渲染误差作为反馈信号。具体地，渲染输入视角的图像，计算与GT的差异作为误差图，将误差图、渲染图像和GT图像的特征拼接送入递归网络，预测高斯参数更新。关键洞察是渲染误差天然适应测试数据分布，减少对训练分布的依赖。无需梯度计算的前馈方式实现迭代优化。其设计目的不是单独增加一个模块，而是用上述表示、监督或渲染约束消除原流程在目标任务中的具体失配。提出前馈递归高斯泼溅模型ReSplat，利用渲染误差作为反馈信号迭代精化3D高斯（无需梯度计算），在16倍下采样空间初始化紧凑高斯集，在多数据集、视角数和分辨率下达到SOTA性能。适用时还应保留论文给出的边界：递归步数增加会线性增加推理时间；初始重建质量依赖DepthSplat的深度预测；递归更新可能收敛到局部最优。
% METHOD: 2512.15508
\paragraph{Off The Grid（子像素关键点高斯分布）（2025.12）}
Off The Grid~\cite{2512_15508} 针对像素网格绑定。它用 patch 内子像素关键点检测决定高斯位置，而不是每个像素固定放一个 primitive。

基于 VGGT 骨干，模型用 U-Net 解码热力图，softmax 得到连续 2D 点坐标；Shannon 熵度量 patch 信息复杂度，并由多密度预测头为不同 patch 分配不同数量高斯。相对 SplatWeaver 的专家路由，它更像从检测角度决定 primitive 放置位置。
关键变化有两层。位置上，卷积热力图和 softmax 坐标让基元中心能落在子像素位置，而不是被像素中心固定；数量上，patch 熵和多密度预测头让简单区域少放高斯、复杂区域多放高斯。它因此同时放松了“在哪里放”和“放几个”的限制。相对 SplatWeaver 的专家路由，Off The Grid 更像把基元生成问题改写成关键点检测与自适应密度估计。
\textbf{实现细节：} 像素对齐方法将高斯基元放置在图像像素网格上，这一刚性约束限制了基元的空间分布灵活性——平坦区域的像素不需要高斯，结构边缘处的像素可能需要多个高斯。Off The Grid（2025.12） 受关键点检测启发，将前馈 3DGS 的基元预测重新设计为亚像素级别的检测问题：解码器学习在每个图像图块（patch）内局部分布基元，取代沿像素网格的刚性排列。同时引入自适应密度机制，根据每个图块的 Shannon 熵分配不同数量的基元——高熵图块（纹理复杂）获得更多基元，低熵图块（平坦区域）获得更少甚至零个基元。整个解码器与预训练三维重建骨干端到端训练，仅使用光度监督而无需任何三维标注。生成的无位姿模型在远少于竞争方法的高斯基元下实现了更高质量的重建，展现了更准确、更高效的基元分配策略。
\textbf{实现补充：} 实现上还要核对输入、输出、参数、特征、初始化、训练损失、正则、渲染或推理接口，以及优化中的剪枝、分裂、采样、编码或解码步骤；这些抓手决定该方法到底改了 3DGS 管线的哪一环。
\textbf{方法信息：} Huawei 参与了这篇论文；代码仓库：\url{https://github.com/ArthurMoreau/OffTheGrid}，核验时完整实现暂未确认。
\textbf{设计动机与边界：} 相对原始 3DGS 以独立高斯和图像重建误差为主的设定，Off The Grid 把问题落在以下可实现的改动上：方法核心是替换传统像素对齐基元分布为自适应Off-The-Grid分布。具体包含三个层次的自适应控制：(1)在图像patch内通过关键点检测技术在子像素级别检测3D高斯基元位置，利用2D关键点作为自监督方式学习最优基元分布，通过卷积热力图提取连续2D点坐标使基元中心脱离网格约束；(2)基于Shannon熵的自适应密度机制，根据patch的信息复杂度分配不同数量的基元；(3)多密度预测头允许每个patch生成不同数量的基元。整体架构基于VGGT-1B骨干网络，去除其pointmap头，使用U-Net卷积模块作为解码器，通过softmax空间维度生成热力图检测基元位置，无需任何3D标注仅用光度监督进行端到端训练。其设计目的不是单独增加一个模块，而是用上述表示、监督或渲染约束消除原流程在目标任务中的具体失配。提出一种新的前馈架构，通过子像素级关键点检测替代像素网格对齐方式来分布3D高斯基元，引入基于Shannon熵的自适应密度机制，结合VGGT骨干网络实现端到端训练，以更少的基元数量达到SOTA新视角合成质量。适用时还应保留论文给出的边界：依赖VGGT骨干的几何预测质量；极端稀疏视图下性能可能受限；超像素级基元检测可能遗漏大面积平坦区域。
% METHOD: 2512.17817
\paragraph{Chorus（多教师 3DGS 场景编码预训练）（2025.12）}
Chorus~\cite{2512_17817} 不直接预测新场景高斯，而是预训练一个 3DGS 场景编码器。它把每个高斯映射为 latent feature，并从多个 2D foundation teachers 蒸馏语义和视觉知识。

SigLIP2 提供语言对齐语义，DINOv3 提供通用视觉特征，PE-Spatial 提供对象感知线索。2D 特征通过渲染权重 uplift 到高斯伪标签，教师特定投影头匹配不同特征空间。相对 TextSplat 的文本引导重建，Chorus 更偏表示预训练和后续任务迁移。
\textbf{实现细节：} Chorus（2025.12） 关注一个被忽视的问题：如何为三维高斯基元学习通用、丰富的特征表示。现有前馈方法通常仅预测高斯的几何和外观参数，而未将通用的视觉理解能力编码到高斯基元中。Chorus 提出多教师预训练框架，从三类互补的二维基础模型蒸馏信号：语言对齐教师（language-aligned teacher，如 CLIP，提供高级语义理解）、通用教师（generalist teacher，提供空间和结构先验）、物体感知教师（object-aware teacher，提供实例级判别能力）。共享的三维编码器通过教师特定的投影头（projector）学习同时捕获从高级语义到细粒度结构的信号。Chorus 在多种任务上评估：开放词汇语义/实例分割、线性/解码器探测、数据高效监督以及基于 LLM 的问答。此外，Chorus 提出渲染-蒸馏适配策略（render-and-distill adaptation）支持域外微调。
\textbf{方法信息：} 代码：\url{https://github.com/gaussianworld/chorus}。
\textbf{设计动机与边界：} 相对原始 3DGS 以独立高斯和图像重建误差为主的设定，Chorus 把问题落在以下可实现的改动上：Chorus框架核心：(1)使用共享3DGS编码器g\_theta将高斯参数映射到潜在逐高斯特征Z；(2)三位2D教师提供互补监督：SigLIP2(语言对齐语义)、DINOv3(通用视觉特征)、PE-Spatial(物体感知线索，结合自对齐与SAM-logit对齐)；(3)每位教师配备轻量级投影头(2层MLP+LayerNorm+GELU)，通过余弦+Smooth-L1匹配损失和可选对比损失进行训练；(4)2D特征通过uplifting(加权平均渲染权重)提升到3D高斯作为伪标签；(5)对教师特征应用PHI-S标准化处理尺度/方差差异；(6)提出render-and-distill适配方法通过图像平面监督加速域外泛化；(7)3DGS感知数据增强辅助预训练。其设计目的不是单独增加一个模块，而是用上述表示、监督或渲染约束消除原流程在目标任务中的具体失配。提出Chorus多教师预训练框架，通过从互补2D基础模型蒸馏学习全面的3DGS场景编码器，使用共享3D编码器和教师特定投影器从语言对齐、通用和物体感知教师中学习，在仅使用39.9倍少训练场景下超越点云基线，提出渲染蒸馏适配方法。适用时还应保留论文给出的边界：编码器质量受限于2D教师的能力；需要预训练的3DGS场景作为输入；大规模预训练计算成本较高。
% METHOD: 2512.18692
\paragraph{EcoSplat（效率可控前馈高斯）（2025.12）}
EcoSplat~\cite{2512_18692} 把目标高斯数量作为条件输入。第一阶段训练标准像素对齐前馈模型；第二阶段做 importance-aware Gaussian finetuning，把目标 primitive 数 $K$ 注入模型，使它学会在不同预算下调整不透明度、协方差和颜色。

基础模型仍用共享 ViT 编码器和交叉注意力解码器，从每个视图 token 预测逐像素高斯中心和属性。IGF 阶段把 $K$ 作为条件信号，让模型在同一输入下学会不同预算的输出策略：不重要基元的不透明度被压低，关键区域的协方差和颜色被保留。推理时，系统按视图重要性分配高斯名额，再保留不透明度最高的 top-K 高斯。相对 TinySplat 的后处理压缩，EcoSplat 在模型内部学习效率-质量折中，因此同一模型可适配不同设备预算。
\textbf{方法信息：} 代码：\url{https://github.com/KAIST-VICLab/EcoSplat}。
\textbf{设计动机与边界：} 相对原始 3DGS 以独立高斯和图像重建误差为主的设定，EcoSplat 把问题落在以下可实现的改动上：EcoSplat采用两阶段优化：(1)像素对齐高斯训练(PGT)：基于现有前馈架构训练像素对齐GS模型，每个视图的token通过共享ViT编码器和交叉注意力解码器处理后，由高斯中心头和参数头预测逐像素高斯基元；(2)重要性感知高斯微调(IGF)：将目标基元数K作为条件信号注入模型，自适应调整高斯参数(不透明度、3D协方差、颜色)，抑制不重要高斯的不透明度，保留最具信息量的基元。推理时根据每个视图的重要性分配逐视图高斯数量，通过保留不透明度最高的K个基元(称为top-K高斯)满足目标约束。关键是K作为可配置条件使同一模型适配不同设备。其设计目的不是单独增加一个模块，而是用上述表示、监督或渲染约束消除原流程在目标任务中的具体失配。提出首个效率可控的前馈3DGS框架EcoSplat，通过两阶段优化(像素对齐训练+重要性感知微调)在推理时自适应预测任意目标基元数量的3D表示，在密集视图设置下严格基元数约束下超越SOTA方法。适用时还应保留论文给出的边界：两阶段训练增加训练复杂度；极低基元数下质量下降明显；基元重要性评估可能不完美。
% METHOD: 2601.03824
\paragraph{IDESplat（迭代深度概率估计）（2026.01）}
IDESplat~\cite{2601_03824} 仍然站在 DepthSplat 这一类 ``先估深度、再反投影成高斯'' 的路线上，但它盯住的是代价体和极线匹配的开销。DepthSplat 把单目深度先验接入前馈 3DGS 后，几何质量明显改善；问题是多视角 warp 若直接保存密集特征，显存会随深度候选和视图数快速增长，而且一次性深度概率估计很难在大视差、遮挡和弱纹理处稳定收敛。IDESplat 的改动是把深度概率估计拆成多轮级联，每轮只保留必要的 warp 索引和概率图，而不是把所有密集 warp 特征都展开存储。

具体做法上，模型先沿候选深度把源视图特征 warp 到目标视图，建立像素之间的极线对应；Warp-Index 机制记录可复用的几何索引，避免重复缓存大块特征。随后 Depth Probability Boosting Unit 将多次级联 warp 得到的极线注意力图以乘法方式融合：高概率深度在多轮中被持续增强，低概率深度被压低。随着轮次推进，深度搜索范围逐步缩小、特征分辨率逐步提高，最后用 softmax 加权得到每个像素的深度期望，再反投影为高斯中心，并由高斯头预测尺度、旋转、不透明度和颜色。相对 DepthSplat，IDESplat 的关键不是换掉 3DGS 表示，而是把深度前端做成可迭代的概率筛选器，从而用更少参数和更低显存得到稳定的像素对齐高斯。\textbf{方法信息：} 代码：\url{https://github.com/cvl-uestc/idesplat}；仓库提供推理、训练和数据配置材料；GitHub 约 42 stars。
\textbf{设计动机与边界：} 相对原始 3DGS 以独立高斯和图像重建误差为主的设定，IDESplat 把问题落在以下可实现的改动上：核心创新：(1)Warp-Index极线注意力——通过warp操作建立源和目标视图之间的像素对应关系，引入warp-index机制减少存储所有密集warp特征的内存开销；(2)Depth Probability Boosting Unit (DPBU)——将多个级联warp操作产生的极线注意力图以乘法方式融合，逐步增强特征相似性度量，识别高似然的深度候选，抑制低概率深度候选；(3)迭代深度估计——通过堆叠多个DPBU构建迭代深度估计过程，在逐步更新深度搜索范围的同时增加特征分辨率，在更精细的尺度上执行warp和相关性计算，逐步精化深度图。最终深度图通过softmax加权平均获得精确Gaussian中心位置。方法在RE10K、ACID、DL3DV和DTU数据集上验证。其设计目的不是单独增加一个模块，而是用上述表示、监督或渲染约束消除原流程在目标任务中的具体失配。提出IDESplat，通过迭代执行warp操作逐步精化深度概率估计，引入Depth Probability Boosting Unit (DPBU)以乘法方式融合级联warp产生的极线注意力图，实现精确的Gaussian均值预测，在仅10.7%参数和70%内存下超越DepthSplat 0.33dB PSNR。适用时还应保留论文给出的边界：迭代warp操作增加了一定的计算开销；在极端大基线多视图场景下可能需要更多迭代次数。
% METHOD: 2603.21304
\paragraph{F4Splat（预测性密度化）（2026.03）}
F4Splat~\cite{2603_21304} 解决前馈高斯均匀分配带来的冗余。模型为每个区域预测 densification score，估计该区域需要的高斯密度，并根据空间复杂度和多视角重叠自适应分配。

用户可以在推理时显式控制最终高斯预算而无需重训练。简单区域分配少，高频或遮挡复杂区域分配多。相对 SplatWeaver 的专家机制，F4Splat 更像一个可调预算的预测性 densification 头。
\textbf{实现细节：} SparseSplat 通过温度系数 $\tau$ 隐式控制密度，但无法直接指定目标基元数量。F4Splat（2026.03） 沿续了自适应密度的思路，但引入致密化分数引导的分配策略（densification-score-guided allocation），将密度分配建模为可学习的逐区域分数预测问题，从而实现显式的基元数控制。具体而言，F4Splat 在标准像素对齐前馈管线的参数预测头中附加一个致密化分数头，为每个空间区域预测一个标量分数，表示该区域所需的高斯密度——分数越高，该区域被分配更多高斯基元。这一设计与 SparseSplat 的关键差异在于密度控制方式：SparseSplat 通过温度系数 $\tau$ 概率性地采样像素，是隐式的密度调节；F4Splat 则显式预测密度需求，并允许用户在推理时指定目标高斯总预算 $K$，模型据此自适应缩放各区域的致密化分数，无需重训练即可在不同基元数量间切换。在实现上，F4Splat 在推理时根据致密化分数对候选位置进行排序，保留分数最高的 $K$ 个位置生成高斯，同时根据空间复杂度和多视角重叠度自适应分配——简单区域和高度重叠视角间分配较少基元，几何复杂且重叠度低的区域获得更多基元。
\textbf{方法信息：} 代码：\url{https://github.com/mlvlab/F4Splat}。
\textbf{设计动机与边界：} 相对原始 3DGS 以独立高斯和图像重建误差为主的设定，F4Splat 把问题落在以下可实现的改动上：F4Splat提出了一种前馈预测性密度化方法。核心创新是密度分数引导的分配策略（densification-score-guided allocation），模型预测每个区域的密度化分数来估计所需的Gaussian密度。与像素到Gaussian或体素到Gaussian的均匀分配不同，该方法根据区域的空间复杂度和多视角重叠程度自适应分配Gaussian。用户可以在推理时显式控制最终Gaussian预算而无需重训练模型。这减少了简单区域的冗余Gaussian，同时最小化重叠视角间的重复Gaussian，产生紧凑且高质量的3D表示。其设计目的不是单独增加一个模块，而是用上述表示、监督或渲染约束消除原流程在目标任务中的具体失配。提出前馈预测性密度化方法F4Splat，解决了前馈3DGS中均匀分配Gaussian导致的冗余问题。引入密度分数引导的分配策略，根据空间复杂度和多视角重叠自适应分配Gaussians，允许在无需重训练的情况下显式控制最终Gaussian预算。在显著减少Gaussian数量的同时实现更优的新视角合成质量。适用时还应保留论文给出的边界：密度分数预测的准确性依赖于训练数据的多样性；对极端场景复杂度的适应性需要进一步验证；与压缩方法的结合潜力未探索。
% METHOD: 2604.01030
\paragraph{Diff3R（不确定性感知可微测试时优化）（2026.04）}
Diff3R~\cite{2604_01030} 试图结合前馈速度和逐场景优化质量。内循环对前馈初始高斯做短暂测试时优化，拟合上下文视角；外循环则用优化后新视角渲染质量更新前馈网络权重。

为了避免展开优化轨迹带来的内存爆炸，Diff3R 使用隐函数定理和 Gauss-Newton 近似求隐梯度，并用矩阵无关 PCG 求解。数据驱动不确定性模型限制欠约束区域的参数变化，防止稀疏视角 TTO 过拟合。相对 ReSplat，它是真正的可微优化层，而不是学习递归更新器。
\textbf{实现细节：} Diff3R（2026.04） 与前述方法的不同在于：它不满足于让前馈网络直接输出最终高斯参数，也不局限于在推理时做若干步梯度精炼，而是将训练过程本身设计为一个双层优化问题——内循环模拟测试时优化（TTO，Test-Time Optimization），外循环据此调整网络，使其学会为 TTO 输出最优的初始化。这一设计源于一个核心观察：前馈网络的单步输出 $\theta_0 = g_\theta(\{I_{\mathrm{ctx}}, P_{\mathrm{ctx}}\})$ 虽然能捕获场景的大致几何和外观，但天然缺少高频细节。而标准 TTO 在稀疏视角下又极易灾难性过拟合——因为没有足够的多视角约束，高斯可以轻易地"作弊"以完美重建输入视角，却严重破坏底层几何。Diff3R 的解决思路是让前馈预测和 TTO 协同训练，而非各自独立运作。双层优化框架。直观上，Diff3R 的训练过程可以理解为一种"模拟考试"：每次训练迭代中，网络先输出初始高斯 $\theta_0$，然后内循环对其执行若干步 3DGS 优化以拟合上下文视角，外循环再根据精炼后 $\theta^{*}$ 在新视角上的渲染质量来更新网络权重。内循环的损失函数为近端正则化的光度损失： 其中 $r(\theta) = f(\theta) - I_{\mathrm{ctx}}$ 为渲染残差，$\Lambda$ 为正则化权重。
\textbf{方法信息：} Huawei 参与了这篇论文。



% METHOD: 2604.07053
\paragraph{AnchorSplat（锚点对齐前馈高斯）（2026.04）}
AnchorSplat~\cite{2604_07053} 也要摆脱像素对齐，但方式是引入 3D anchors。它先用多视角立体模块或 RGB-D 点云预测可靠几何先验，再把多视图特征聚合到这些 3D 锚点上。

锚点预测器可来自 MapAnything、MVSAnywhere 等多视角立体模块，也可以来自稀疏点云、体素或 RGB-D 点云。CNN 编码多视图外观后，Transformer 解码器围绕这些锚点预测位置、尺度、不透明度和颜色；一个即插即用高斯精炼器再用少数前向步骤修正中间高斯属性，无需重训完整模型。相对 VolSplat 的规则体素，AnchorSplat 的支撑是稀疏但语义明确的 3D 几何锚点，因此更容易减少冗余高斯和跨视角浮点。
\textbf{实现细节：} 与 GlobalSplat 的全局 token 表示不同，AnchorSplat（2026.04） 提出锚点对齐的高斯表示，与 TokenGS 的可学习 token 方案不同，AnchorSplat 以外部提供的三维几何先验（稀疏点云、体素或 RGB-D 点云）作为锚点，直接在三维空间中构建高斯表示。具体而言，AnchorSplat 将多视角图像特征反投影到锚点位置形成锚定特征体（anchored feature volume），再由高斯解码器从中预测高斯参数——这一流程使高斯表示完全独立于图像分辨率和视角数量。此外，AnchorSplat 引入即插即用的 Gaussian Refiner 模块，通过数次前向传播精调中间高斯的属性（位置、尺度、颜色），提升几何一致性而无需重训练整个模型。
\textbf{实现补充：} 实现上还要核对输入、输出、参数、特征、初始化、训练损失、正则、渲染或推理接口，以及优化中的剪枝、分裂、采样、编码或解码步骤；这些抓手决定该方法到底改了 3DGS 管线的哪一环。
\textbf{设计动机与边界：} 相对原始 3DGS 以独立高斯和图像重建误差为主的设定，AnchorSplat 把问题落在以下可实现的改动上：AnchorSplat包含三个核心组件：(1) 锚点预测器：利用预训练的多视角立体模块（如MapAnything、MVSAnywhere等）从输入图像预测可靠的3D几何先验和深度图作为空间锚点；(2) 高斯解码器：CNN特征提取器编码多视角信息后，通过Transformer解码器预测锚点对齐的高斯，确保与3D先验一致；(3) 高斯精炼器：即插即用模块，仅需几次前向传播即可调整高斯属性（位置、尺度、不透明度、颜色），无需重新训练完整模型。关键设计是将多视角特征聚合到一致的3D锚点后再预测高斯，而非从2D像素直接映射，从而消除浮子和跨视角不一致性。其设计目的不是单独增加一个模块，而是用上述表示、监督或渲染约束消除原流程在目标任务中的具体失配。提出AnchorSplat，首个锚点对齐的前馈3DGS框架，直接在3D空间中表示场景而非绑定到2D像素。利用3D几何先验（稀疏点云、体素或RGB-D点云）引导高斯预测，大幅减少所需高斯数量，提升计算效率和重建保真度。引入即插即用的高斯精炼器，仅需几次前向传播即可调整中间高斯。在ScanNet++ v2 NVS基准上达到SOTA。适用时还应保留论文给出的边界：依赖外部MVS模块的质量；锚点数量和分布影响最终重建质量；可能在大规模开放场景中面临扩展性挑战。
\textbf{方法信息：} Huawei 参与了这篇论文；代码：\url{https://github.com/Zhang-Xiaoxue/AnchorSplat}。



% METHOD: 2604.08370
\paragraph{SurfelSplat（Nyquist 引导的前馈 surfel）（2026.04）}
SurfelSplat~\cite{2604_08370} 从采样理论解释前馈高斯 surfel 几何失败。像素对齐 2D 高斯的空间频率可能超过采样率，导致法线估计不准、surfel 倾向平行图像平面。

方法先根据焦距和深度计算每个 surfel 的空间采样率，再用 Nyquist 引导的低通滤波器调整 surfel 几何，使其频率内容在可采样范围内。之后把适配 surfel 投影到所有输入视角聚合特征，并由融合网络回归表面对齐的 2D 高斯。相对 SparseSurf，它是前馈而非逐场景优化。
\textbf{实现细节：} SurfelSplat（2026.04） 从信号处理的基本原理出发，为前馈 3DGS 的几何精度问题提供了一个新视角。作者观察到，像素对齐高斯基元的空间频率天然超过 Nyquist 采样率——即高斯基元在图像平面上的投影密度过高，相邻基元之间的空间间隔不足以准确恢复每个基元的几何属性（法向量、曲率等）。这导致虽然渲染的 RGB 图像看起来合理，但提取的几何表面却充满噪声和断裂。为解决这一问题，SurfelSplat 提出基于采样率引导的低通滤波策略：首先根据输入分辨率和基元密度计算局部 Nyquist 频率，然后对高斯基元的几何属性施加空间低通滤波，滤除超出采样率的高频分量。滤波后的基元被投影到所有输入视角以获取跨视角特征相关性，通过专门设计的特征融合网络回归具有精确几何的高斯 surfel 表示。
\textbf{方法信息：} 代码：\url{https://github.com/simon-dcs/surfel_splat}。
\textbf{设计动机与边界：} 相对原始 3DGS 以独立高斯和图像重建误差为主的设定，SurfelSplat 把问题落在以下可实现的改动上：SurfelSplat的核心贡献是首次从Nyquist采样定理角度分析前馈高斯surfel的几何恢复问题。关键发现：像素对齐的2D高斯surfel的空间频率超过Nyquist采样率，导致法线无法精确恢复（surfel趋向平行于图像平面而非对齐实际表面）。解决方案包含三个步骤：(1) 计算每个surfel的空间采样率：基于相机焦距和深度推导surfel覆盖的图像区域；(2) Nyquist引导的surfel适配：用空间采样率引导的低通滤波器调整surfel几何形式，确保其频率内容在采样能力范围内；(3) 跨视角特征聚合：将适配后的surfel投影到所有输入视角获取特征关联，通过专门设计的融合网络回归表面对齐的2D高斯辐射场。整个流程可在1秒内完成。其设计目的不是单独增加一个模块，而是用上述表示、监督或渲染约束消除原流程在目标任务中的具体失配。提出SurfelSplat，首个基于Nyquist采样定理分析前馈高斯surfel几何恢复失败原因的框架。发现像素对齐基元的空间频率超过Nyquist采样率导致几何属性无法精确恢复。提出跨视角特征聚合模块：先用空间采样率引导的低通滤波器适配surfel几何形式，再投影到所有输入视角获取跨视角特征关联，通过特征融合网络回归精确几何的surfel。1秒内预测，比逐场景优化快100倍。适用时还应保留论文给出的边界：主要在DTU等受控数据集上评估；对极端稀疏视角（如单视角）的处理可能不足；低通滤波可能损失高频细节。
% METHOD: 2604.09903
\paragraph{PointSplat（几何驱动剪枝与 Transformer 精炼）（2026.04）}
PointSplat~\cite{2604_09903} 处理已有高斯模型的高效剪枝和精炼，但不依赖 2D 图像。它只看高斯的 3D 属性，如体积、不透明度、位置和尺度，计算剪枝分数。

剪枝分数由归一化不透明度和体积等属性组合而成，用来在空间覆盖和表达贡献之间做平衡；这避免了重新遍历图像射线或计算像素交叉。保留下来的高斯进入双分支编码器：几何分支处理位置和尺度，外观分支处理球谐颜色，Transformer 再统一精炼属性。相对 SplatFormer 的 OOD 修复，PointSplat 更强调纯 3D 几何驱动的压缩与效率。
\textbf{实现细节：} 上述三种方法（GlobalSplat、AnchorSplat、VolSplat）分别从全局 token、几何锚点和体素空间三个维度突破了像素对齐的表示瓶颈，但都聚焦于前馈重建阶段。PointSplat（2026.04） 则从另一个角度切入——结合三维几何驱动的剪枝和 Transformer 精炼，仅依赖高斯基元的三维属性进行重要性排序，无需二维图像参与剪枝阶段。与 SparseSplat 在前馈管线中自适应采样不同，PointSplat 面向的是后处理场景——给定一个已生成的过参数化 3DGS 模型，如何高效压缩。其核心贡献有两个：(1) 几何驱动的剪枝分数——仅基于高斯的不透明度和体积（三维空间覆盖）计算重要性分数，平衡空间覆盖与信息密度，避免了传统方法需要光线-像素交点计算的昂贵开销；(2) 双分支编码器（Dual-Branch Encoder）——将几何属性（位置、尺度、朝向）与高维外观特征（球谐系数）显式分离，通过独立分支分别处理并重新加权，避免高维外观特征对低维几何属性的支配（feature imbalance）。剪枝后的高斯经 Transformer 精炼模块全局协调参数，修复因剪枝引入的模糊纹理和断裂几何边缘。
\textbf{方法信息：} 代码：\url{https://github.com/anhthuan1999/pointsplat}。
\textbf{设计动机与边界：} 相对原始 3DGS 以独立高斯和图像重建误差为主的设定，PointSplat 把问题落在以下可实现的改动上：方法包含两个核心组件：(1)几何驱动剪枝策略：仅利用高斯的内在3D属性计算剪枝分数，平衡空间覆盖和不透明度，通过z-score归一化的不透明度和体积的加权和作为分数，选择紧凑但信息丰富的子集，完全无需2D图像和射线-像素交叉计算；(2)双分支编码器：分别处理几何特征(位置、尺度)和外观特征(球谐系数)，通过重新加权避免两类特征不平衡。使用Transformer架构进行精炼，使保留的高斯在属性上得到优化。其设计目的不是单独增加一个模块，而是用上述表示、监督或渲染约束消除原流程在目标任务中的具体失配。提出PointSplat，首个纯3D几何驱动的剪枝-精炼框架。仅基于高斯的3D属性(体积和不透明度)计算剪枝分数，无需2D图像和逐场景优化。设计双分支编码器分离并重新加权几何和外观特征避免特征不平衡。在ScanNet++和Replica上跨多种稀疏度级别保持竞争力渲染质量的同时显著提升效率。适用时还应保留论文给出的边界：极端剪枝率下渲染质量可能下降，依赖初始3DGS模型的质量，对纹理丰富场景的保持能力有待进一步验证。
% METHOD: 2604.15239
\paragraph{TokenGS（可学习高斯 token）（2026.04）}
TokenGS~\cite{2604_15239} 将高斯预测从像素网格解耦。它不沿每个像素射线回归深度，而是引入固定数量的 learnable Gaussian tokens，通过交叉注意力从图像特征中吸收信息，再直接预测 3D 坐标和其他高斯参数。

高斯数量变成与输入分辨率和视角数无关的超参数。visibility-aware 正则解决所有相机视锥外高斯无梯度的问题；token 空间还支持自监督测试时优化。相对 C3G 和 GlobalSplat，它的重点是把 primitive 本身作为 token。
\textbf{实现细节：} 前述自适应密度方法虽然缓解了冗余问题，但高斯位置仍受输入像素或射线约束。TokenGS（2026.04） 完全脱离了像素对齐这一路线——无论是像素对齐的 MVSGaussian/DepthSplat（相应公式），还是自适应密度的 SparseSplat（相应公式），TokenGS 都抛弃了"高斯位置由输入像素/射线决定"这一根本假设，使高斯的位置和数量都与输入图像解耦。基线与问题：像素对齐前馈 3DGS 存在三个根本性限制。第一，射线约束——将高斯均值参数化为沿相机射线的深度值，使模型无法内部修正噪声位姿和多视角不一致性，且难以扩展到动态场景（需要随时间 warp 点/像素）。第二，数量耦合——高斯数量与输入分辨率和视角数绑定，例如 32 张 $512 \times 512$ 图像产生超过 800 万个粒子，且将同一输入重复 $N$ 次就产生 $N$ 倍粒子，无论场景本身复杂度如何。第三，测试时微调破坏先验——直接优化高斯参数容易退化网络学到的先验，尤其在少视角场景中。核心改动：TokenGS 针对上述三个限制做了对应的三个关键改变。
\textbf{方法信息：} NVIDIA 参与了这篇论文。
\textbf{设计动机与边界：} 相对原始 3DGS 以独立高斯和图像重建误差为主的设定，TokenGS 把问题落在以下可实现的改动上：TokenGS的核心创新包括：(1) 直接回归3D高斯均值坐标而非沿相机射线的深度，通过自监督渲染损失训练，无需显式3D监督；(2) 编码器-解码器架构中，可学习高斯token通过交叉注意力获取图像特征并预测3DGS参数，高斯数量成为与输入无关的超参数；(3) 引入visibility-aware正则化解决零梯度问题（处于所有相机视锥外的高斯无梯度）；(4) token空间的自监督测试时优化保留模型全局先验。方法自然支持静态-动态分解和场景流恢复。其设计目的不是单独增加一个模块，而是用上述表示、监督或渲染约束消除原流程在目标任务中的具体失配。提出TokenGS，将高斯均值预测从像素对齐的深度回归改为直接回归3D坐标，引入编码器-解码器架构配合可学习高斯token，解除预测原语数量与输入图像分辨率和视角数量的绑定关系。支持token空间的高效测试时优化而不退化学习先验，在静态和动态场景上均达到SOTA前馈重建性能。适用时还应保留论文给出的边界：token数量作为超参数需手动设置；对极端稀疏视角可能仍有挑战；动态场景处理需要额外的时间建模。
% METHOD: 2604.15284
\paragraph{GlobalSplat（全局场景 token）（2026.04）}
GlobalSplat~\cite{2604_15284} 采用“先对齐后解码”原则。它先把所有输入视角融合成全局对齐的 latent scene tokens，再从这些全局 token 解码 3D 高斯。

输入 RGB patch token 注入 Plucker ray 和全局相机编码；双分支注意力架构迭代细化场景 token，粗到细容量课程逐步增加解码能力，避免表示膨胀。相对 TokenGS 的高斯 token，GlobalSplat 的 token 更像场景记忆，数量恒定且与输入视角无关。
它先用相似变换把相机和场景归一到规范坐标系，再把 patch 化 RGB token、Plucker 射线嵌入和全局相机编码组成输入上下文。潜变量场景 token 通过交叉注意力反复读取多视角证据，并在由粗到精的课程中逐步增加解码容量。这里的关键约束是 token 数量恒定：输入视角越多，只是提供更多上下文，不会线性增加高斯候选。相对 TokenGS 的“每个 token 直接成为 primitive”，GlobalSplat 更像用固定容量场景记忆先汇总全局结构，再解码高斯。
\textbf{实现细节：} TokenGS 通过可学习 token 直接回归全局三维坐标，实现了位置与输入的解耦。GlobalSplat（2026.04） 遵循与 TokenGS 不同的"先对齐，后解码"（align first, decode later）原则，通过双分支迭代注意力架构在潜在空间中显式构建全局场景表示，再从中解码三维几何。TokenGS 通过可学习 token 直接回归全局三维坐标，而解码前的对齐仍依赖 Transformer 的隐式学习。GlobalSplat 则将"先对齐"显式化为一个独立的潜在场景表示构建阶段——先用双分支迭代注意力架构（显式解耦几何分支和外观分支，各自拥有独立的查询 $Q_G$、$Q_A$）将所有输入视角的特征融合到一组固定数量的全局潜在场景 token 中，跨视角对应关系在这一阶段通过交叉注意力被显式解决。只有在全局场景表示构建完成后，才开始解码显式三维几何。
\textbf{方法信息：} 代码：\url{https://github.com/R-Itk/globalsplat}。
\textbf{设计动机与边界：} 相对原始 3DGS 以独立高斯和图像重建误差为主的设定，GlobalSplat 把问题落在以下可实现的改动上：GlobalSplat采用'先对齐后解码'范式：(1) 场景归一化：将场景映射到规范坐标系，使用相似变换使相机姿态一致；(2) 输入上下文构建：提取patch化的RGB token，注入Plucker射线嵌入和全局相机编码；(3) 双分支注意力架构：迭代细化潜变量场景token，通过交叉注意力融合多视角信息；(4) 由粗到精容量课程：逐步增加解码容量，严格防止表示膨胀。场景token数量恒定，与输入视角数量无关。其设计目的不是单独增加一个模块，而是用上述表示、监督或渲染约束消除原流程在目标任务中的具体失配。提出GlobalSplat，基于'先对齐后解码'原则，先融合所有输入视角到全局对齐的潜变量场景表示，再解码为3D高斯。使用双分支注意力架构和由粗到精的训练课程，仅需16K高斯即可在RealEstate10K和ACID上实现竞争力性能，模型仅4MB，推理时间78ms以内。适用时还应保留论文给出的边界：全局token可能丢失细节信息；由粗到精训练需要精心设计课程；对非常复杂场景可能需要更多token。
% METHOD: 2605.07287
\paragraph{SplatWeaver（自适应高斯基数专家）（2026.05）}
SplatWeaver~\cite{2605_07287} 回到前馈 3DGS 的一个基础假设：很多像素对齐方法默认每个像素生成固定数量的高斯。这个规则实现简单，但平滑墙面会浪费基元，枝叶、细杆、纹理边缘又不够用。SplatWeaver 的改变是把 ``生成几个高斯'' 也交给网络学习，用 cardinality Gaussian experts 做区域自适应分配。

流程上，模型先用 DINOv2 和多视图几何 Transformer 提取跨视图特征，再由 DPT 解码器得到像素级表示。随后，高频先验模块从离散小波域提取纹理和细结构信号，交给路由器判断该位置应该调用哪个基数专家。第 $m$ 个专家专门生成 $m$ 个高斯，$m$ 从 0 到上限 $M$；平滑区域可以选 0 或少量高斯，复杂区域可以选更多高斯。路由正则用于避免专家选择震荡，并让不同专家真正承担不同复杂度区域。

相对 FreeSplat 的跨视角冗余融合，SplatWeaver 更早一步在生成阶段控制 primitive 数量；相对 EcoSplat 一类全局预算方法，它把预算分配细化到像素或局部区域。\textbf{方法信息：} Huawei 参与了这篇论文；公开仓库：\url{https://github.com/yecongwan/splatweaver}；GitHub 约 128 stars。需要注意的是，README 显示仓库提供推理和训练/评测 demo，但完整训练实现仍列为待办，因此本书把它作为公开仓库和高关注度方法记录，不把它写成完整训练代码已发布。
\textbf{设计动机与边界：} 相对原始 3DGS 以独立高斯和图像重建误差为主的设定，SplatWeaver 把问题落在以下可实现的改动上：核心创新在于自适应原语分配：(1)基数Gaussian专家：设计M+1个专家(产生0到M个Gaussian)，每个专家专精于特定数量的原语生成；(2)像素级路由：为每个像素特征分配最合适的基数专家，实现空间自适应分配；(3)高频先验引导：从离散小波域提取高频先验，引导路由为细结构和纹理区域分配更多Gaussian；(4)路由正则化：稳定专家选择，促进复杂度感知分配。整体流程：DINOv2 patch化→多视图几何Transformer→DPT解码器→高频先验→专家路由→Gaussian预测。支持未标定输入，联合估计相机位姿。其设计目的不是单独增加一个模块，而是用上述表示、监督或渲染约束消除原流程在目标任务中的具体失配。提出基数Gaussian专家机制，每个专家负责产生特定数量的原语(0到M)，通过像素级路由自适应分配。引入高频先验引导模块和路由正则化。在多种场景上一致超越SOTA方法，使用更少Gaussian实现更忠实的新视角渲染。适用时还应保留论文给出的边界：路由机制的额外计算开销；专家数量M需预设上限；对极端复杂场景的自适应能力受限。
% METHOD: 2605.10239
\paragraph{AdaptSplat（VFM 高频适配器）（2026.05）}
AdaptSplat~\cite{2605_10239} 认为复杂前馈组件不一定必要，关键是补偿视觉基础模型深层特征的低通倾向。它用轻量 Feature Prior Adapter 从 DINOv3-ConvNeXt 浅层特征中提取方向感知高频结构。

FPA 通过 2D 小波分解得到低频和高频子带，再用高频位置编码和自适应残差调制注入 DPT 解码器，防止高斯尺度均匀收敛为球形。相对 HiSplat 的多尺度高斯补偿，AdaptSplat 更像骨干适配器。
实现上，DINOv3-ConvNeXt 骨干完全解冻，标准多视图 Transformer 负责跨视角交互，DPT 解码器回归空间结构。FPA 只在关键位置补上高频几何先验：浅层特征经 2D 离散小波变换分成 LL/LH/HL/HH 子带，方向性高频再通过位置编码和残差调制进入解码器。这样，模型在预测高斯尺度时能感知边缘方向和细结构，不容易把所有尺度收敛成缺少方向性的球形。
\textbf{方法信息：} 代码：\url{https://github.com/xmw666/adaptsplat}。
\textbf{设计动机与边界：} 相对原始 3DGS 以独立高斯和图像重建误差为主的设定，AdaptSplat 把问题落在以下可实现的改动上：使用DINOv3-ConvNeXt作为骨干（完全解冻端到端训练），标准多视图Transformer做跨视角交互，标准DPT解码器做空间回归。关键创新是FPA适配器：用2D离散小波变换(DWT)分解浅层特征，提取LL/LH/HL/HH子带中的方向感知高频结构先验。这些先验通过高频位置编码和自适应残差调制融入解码器，防止缩放系数均匀收敛导致的各向同性（球形）退化。其设计目的不是单独增加一个模块，而是用上述表示、监督或渲染约束消除原流程在目标任务中的具体失配。提出AdaptSplat，证明仅需1.5M参数的轻量适配器即可超越复杂组件设计。FPA从VFM浅层特征提取方向感知高频结构先验，通过高频位置编码和自适应残差调制融入通用管线，有效补偿深层网络的低通滤波效应，避免高斯的各向同性退化。适用时还应保留论文给出的边界：依赖VFM质量；对极端视角变化仍有局限。
% METHOD: 2605.13093
\paragraph{RoSplat（视角数鲁棒像素级前馈）（2026.05）}
RoSplat~\cite{2605_13093} 处理输入视角数量变化和高分辨率渲染。像素级前馈方法在视角数增加时，alpha blending 数值可能导致结果过亮；RoSplat 提出视角数自适应归一化，使渲染结果不随输入数量漂移。

同时，它补强高斯尺度监督，用多视角深度一致性约束尺度估计。相对 FreeSplat 的长序列聚合，RoSplat 更专注于“同一模型在不同输入视角数量下稳定工作”。
它的自适应归一化直接作用在 $\alpha$-混合数值上，抵消输入视角数量变化带来的亮度和透射率漂移；尺度感知监督则用多视角深度一致性约束高斯大小，避免模型只学到中心位置而忽略 footprint。相对 FreeSplat 的长序列聚合，RoSplat 更专注于“同一模型在不同输入视角数量下稳定工作”。
\textbf{实现细节：} 像素对齐前馈方法的一个被忽视的缺陷是：当推理时的输入视角数量与训练时不一致，渲染图像会出现系统性过亮。RoSplat（2026.05） 分析了这一现象的根源：像素对齐高斯按深度排序执行 $\alpha$-混合时，重叠高斯的数量随视角数线性增长，每个额外视角的高斯贡献额外的颜色累积，导致亮度偏移。RoSplat 的解决方案简洁而有效。Alpha 归一化：在标准 $\alpha$-混合渲染后，同时渲染累积 alpha 值 $\alpha(\mathbf{x}) = \sum_i \alpha_i^{\mathrm{clip}} G_i^{2D}(\mathbf{x}) \prod_{j<i}(1 - \alpha_j^{\mathrm{clip}} G_j^{2D}(\mathbf{x}))$，以最终颜色除以 alpha 值 $C_{\mathrm{norm}}(\mathbf{x}) = C(\mathbf{x}) / \alpha(\mathbf{x})$——直觉上，视角越多 alpha 累积越高，归一化后颜色恢复到物理正确的亮度水平。这一策略不改变训练过程，仅在推理时应用。三维采样正则化器解决高分辨率渲染中的孔洞伪影：高斯尺度估计不准确时，在高分辨率下放大渲染会暴露高斯之间的间隙。RoSplat 引入三维空间中的辅助采样约束，鼓励高斯尺度预测的准确性——通过在训练中对高斯三维位置施加采样一致性约束，使高斯更紧密地覆盖表面，减少高分辨率下的间隙。
\textbf{技术实现补充：} 像素对齐前馈方法的一个被忽视的缺陷是：当推理时的输入视角数量与训练时不一致，渲染图像会出现系统性过亮。RoSplat（2026.05） 分析了这一现象的根源：像素对齐高斯按深度排序执行 $\alpha$-混合时，重叠高斯的数量随视角数线性增长，每个额外视角的高斯贡献额外的颜色累积，导致亮度偏移。RoSplat 的解决方案简洁而有效。Alpha 归一化：在标准 $\alpha$-混合渲染后，同时渲染累积 alpha 值 $\alpha(\mathbf{x}) = \sum_i \alpha_i^{\mathrm{clip}} G_i^{2D}(\mathbf{x}) \prod_{j<i}(1 - \alpha_j^{\mathrm{clip}} G_j^{2D}(\mathbf{x}))$，以最终颜色除以 alpha 值 $C_{\mathrm{norm}}(\mathbf{x}) = C(\mathbf{x}) / \alpha(\mathbf{x})$——直觉上，视角越多 alpha 累积越高，归一化后颜色恢复到物理正确的亮度水平。这一策略不改变训练过程，仅在推理时应用。三维采样正则化器解决高分辨率渲染中的孔洞伪影：高斯尺度估计不准确时，在高分辨率下放大渲染会暴露高斯之间的间隙。RoSplat 引入三维空间中的辅助采样约束，鼓励高斯尺度预测的准确性——通过在训练中对高斯三维位置施加采样一致性约束，使高斯更紧密地覆盖表面，减少高分辨率下的间隙。
\textbf{方法信息：} NVIDIA 参与了这篇论文。
\textbf{设计动机与边界：} 相对原始 3DGS 以独立高斯和图像重建误差为主的设定，RoSplat 把问题落在以下可实现的改动上：两个互补组件：(1)视角数鲁棒渲染：分析alpha-blending在视角数变化时的数值行为，提出自适应归一化机制，使渲染结果不受输入视角数影响；(2)尺度感知监督：现有方法对高斯尺度的监督不足，提出利用多视角几何约束（深度一致性）增强尺度估计精度。其设计目的不是单独增加一个模块，而是用上述表示、监督或渲染约束消除原流程在目标任务中的具体失配。解决前馈像素级3DGS的两个关键问题：(1)输入视角数变化导致渲染过亮；(2)高斯尺度估计监督不足。提出视角数自适应渲染机制和尺度感知监督，提升跨设置鲁棒性。适用时还应保留论文给出的边界：对极少视角（如2视图）的改善有限。

\section{无标定与自由视角}

Pose-free 前馈方法进一步放宽条件：输入图像可能没有外参、没有内参，甚至视角重叠小、外观变化大。这里的锚点是 PF3plat、FreeSplatter、NoPoSplat、Splatt3R 和 AnySplat：它们分别用预训练几何模型、规范空间、MASt3R 扩展、VGGT 蒸馏或统一相机 token 解决位姿与高斯预测的耦合。
% METHOD: 2408.13912
\paragraph{Splatt3R（MASt3R 扩展到高斯）（2024.08）}
Splatt3R~\cite{2408_13912} 是 pose-free 图像对高斯预测的早期代表。它基于 MASt3R，原模型能从无标定图像对预测 3D 点云；Splatt3R 增加第三分支，预测不透明度、协方差和颜色等 Gaussian 属性。

输入两张未标定图像时，模型单次前向输出两个视角的像素对齐 3D Gaussians，可直接渲染新视角。训练分两阶段：先用 3D 点云几何损失稳定结构，再用新视角合成损失微调外观；外推视角使用空间距离相关的 loss mask，避免强外推区域的错误监督主导训练。相对 FreeSplatter，Splatt3R 更直接地继承 MASt3R 几何能力，核心改动就是把“点云基础模型”扩展为“点云 + 高斯属性”的可渲染前馈模型。
\textbf{实现细节：} Splatt3R（2024.08） 走了一条与前述方法不同的路线——不从头训练前馈 3DGS 网络，而是直接在三维几何基础模型 MASt3R 的输出上追加高斯属性。其核心洞察在于：MASt3R 预测的像素对齐三维点与通用 3DGS 方法预测的像素对齐高斯基元在结构上高度相似——二者都基于跨注意力 Transformer 架构提取多视角信息，并为每个像素预测三维位置。因此，只需在 MASt3R 的点云输出上追加高斯基元所需的额外属性（协方差、不透明度、颜色），即可将其转化为完整的三维高斯表示，无需相机标定。在具体实现上，Splatt3R 在冻结的 MASt3R 模型上新增第三条预测头——Gaussian head，与原有的点云预测头和特征匹配头并行运行。该头同样基于 DPT 架构，为每个像素预测：(1) 协方差，参数化为四元数旋转 $q \in \mathbb{R}^{}4$ 和缩放 $s \in \mathbb{R}^{}3$；(2) 球谐颜色参数 $S \in \mathbb{R}^{3 \times d}$；(3) 不透明度 $\alpha \in \mathbb{R}$；(4) 位置偏移 $\delta \in \mathbb{R}^{}3$。高斯均值参数化为 $\mu = x + \delta$，其中 $x$ 是 MASt3R 预测的三维点。所有参数均在第一帧的相机坐标系中表示，避免使用真值变换在不同参考帧间转换参数。各参数类型使用对应的激活函数：四元数归一化、缩放和偏移指数激活、不透明度 sigmoid 激活。为辅助高频颜色学习，模型预测像素颜色与高斯颜色之间的残差。最终的三维高斯集合是两幅输入图像预测的高斯的并集。关于训练策略，Splatt3R 的设计体现了精心的考量与平衡。MASt3R 的点云预测为每个像素提供了从损失函数到正确三维位置的直接梯度路径，其具体形式为$L_\mathrm{pts} = \sum_{v} \sum_i C_i^v \left(\frac{1}{z}\|X_i^v - \hat{X}_i^v\| - \lambda \log(C_i^v)\right)$，这使得高斯训练不会陷入局部最优（原始 3DGS 优化的核心痛点）。训练采用两阶段方案：第一阶段使用 MASt3R 预训练的点云几何损失提供稳定的三维初始化；第二阶段冻结 MASt3R 的其余参数，仅训练 Gaussian head，使用掩码重建损失优化新视角合成质量： 其中 $M$ 为损失掩码——通过将目标图像的像素反投影并重投影到上下文图像中，检查渲染深度是否与真值深度匹配来识别可见像素，仅对在至少一幅上下文图像中可见的区域计算渲染损失。
\textbf{设计动机与边界：} 相对原始 3DGS 以独立高斯和图像重建误差为主的设定，Splatt3R 把问题落在以下可实现的改动上：Splatt3R基于MASt3R(3D几何重建基础模型)扩展。MASt3R原本只重建3D点云，Splatt3R为其添加第三分支预测额外Gaussian属性(不透明度、协方差、颜色)。训练策略分两阶段：(1)先用3D点云几何损失训练，学习稳定的3D结构；(2)再用新视角合成损失微调外观。提出新型损失掩码策略，对外推视角使用空间距离相关的掩码权重，这对强外推视角的性能至关重要。推理时，两个无标定图像输入模型，单次前向传播输出两个视角的像素对齐3D Gaussian，可直接渲染新视角。其设计目的不是单独增加一个模块，而是用上述表示、监督或渲染约束消除原流程在目标任务中的具体失配。提出Splatt3R，首个从无标定图像对直接预测3D Gaussian的前馈方法，基于MASt3R扩展预测Gaussian属性，通过两阶段训练避免局部最优，实现4FPS的实时重建。适用时还应保留论文给出的边界：仅支持两个视角输入；视角依赖的颜色(SH)建模仍有挑战；重建分辨率受限于512x512。
\textbf{方法信息：} 代码：\url{https://github.com/harim-54/harimsplatt3r}；引用很多，Semantic Scholar 约 170 次。



% METHOD: 2410.22128
\paragraph{PF3plat（无位姿单次前馈 3DGS）（2024.10）}
\textbf{锚点细化：} PF3plat：针对稀疏视角、噪声位姿和低重叠等非理想条件下训练不稳定的问题，用``几何置信度引导的三阶段精炼''替代 NoPoSplat 的直接训练，核心思想是第一阶段用预训练单目深度和视觉对应模型实现跨视角高斯粗对齐，第二阶段精炼深度和位姿估计，第三阶段计算几何置信度分数条件化高斯参数预测——对低置信度区域施加更强正则化抑制梯度更新，对高置信度区域允许更灵活调整，有效区分几何可靠与不可靠区域。
PF3plat~\cite{2410_22128} 是 pose-free 前馈 3DGS 的核心锚点之一。它指出无位姿输入下最先出错的不是颜色头，而是高斯中心：像素对齐方法需要知道每个像素沿哪条射线、在多远处落到 3D 空间；深度或相机稍偏，来自不同视角的高斯就不能对齐，后续协方差、不透明度和颜色预测都会受到污染。

PF3plat 用三阶段几何管线稳定这个问题。第一阶段调用预训练单目深度模型 UniDepth 和视觉对应模型 MASt3R，得到粗深度和粗相对位姿，把每个视角的像素高斯先放到大致一致的空间中。第二阶段用轻量 Transformer 模块对粗估计做残差精修：深度分支预测深度校正，位姿分支预测相对位姿偏移。第三阶段根据精修后的深度和跨视角对应计算几何置信度，并把这个置信度作为条件输入到高斯参数预测头；高置信区域可以生成更有表达力的高斯，低置信区域则保守处理，避免错误中心带来大面积伪影。这个设计也让训练信号更干净：网络不用在所有像素上同等相信粗几何，而是按可靠性分配更新强度，尤其适合视角重叠少、遮挡多的输入组合。

与 NoPoSplat 这类更依赖端到端光度训练的规范空间方法相比，PF3plat 更像 ``foundation model 粗对齐 + 轻量可学习修正 + 置信度调制''。它的关键不是取消相机估计，而是承认相机和深度估计不可避免有误差，再把误差可靠性显式传给高斯属性预测。几何置信度在这里不只是一个后验 mask：它参与调制不透明度、协方差和颜色头，让网络知道哪些中心可以承载锐利细节，哪些中心更适合降低权重以避免漂浮。这样，PF3plat 把 pose-free 的不确定性从隐式训练难题变成了属性预测时可利用的输入变量。\textbf{方法信息：} Microsoft Research 参与了这篇论文；代码：\url{https://github.com/cvlab-kaist/pf3plat}；仓库提供训练和评测脚本；GitHub 约 245 stars，Semantic Scholar 约 63 次引用。

\textbf{实现细节：} PF3plat（2024.10） 在 NoPoSplat 的坐标规范化框架上，针对稀疏视角、噪声位姿和低重叠等非理想条件下的训练不稳定性问题，提出了三阶段精炼策略。第一阶段，利用预训练单目深度估计和视觉对应模型实现跨视角高斯的粗对齐——为不同视角预测的高斯提供初始的空间一致性。第二阶段，引入轻量可学习模块从粗对齐出发精炼深度和位姿估计，提高三维重建和新视角合成质量。第三阶段，利用精炼后的估计计算几何置信度分数——评估每个高斯中心的可靠性——并据此条件化高斯参数预测：对低置信度区域的高斯施加更强的正则化约束（抑制其在训练中的梯度更新幅度），对高置信度区域则允许更灵活的参数调整。这一置信度感知策略有效区分了几何可靠与不可靠的区域，使模型在弱几何约束下仍能产生鲁棒的重建。在大规模真实世界数据集的广泛评估中，PF3plat 在所有基准上均达到最佳性能。
\textbf{设计动机与边界：} 相对原始 3DGS 以独立高斯和图像重建误差为主的设定，PF3plat 把问题落在以下可实现的改动上：方法解决像素对齐3DGS在宽松假设下的挑战：1)粗对齐：使用预训练单目深度估计（UniDepth）和视觉对应模型（MASt3R）获取初始深度和相机位姿估计，实现3D高斯的粗略空间对齐。2)精修模块：轻量级可学习Transformer模块，以粗估计为输入精修深度和位姿。深度精修模块预测残差深度校正，位姿精修模块预测相对位姿偏移。3)几何置信度评分：基于精修后的深度和对应关系计算每个3D高斯中心的可靠性评分，用于条件化高斯参数（不透明度/协方差/颜色）的预测。高置信度区域生成高质量高斯，低置信度区域保守处理避免伪影。整个框架单次前馈，无需场景特定优化。其设计目的不是单独增加一个模块，而是用上述表示、监督或渲染约束消除原流程在目标任务中的具体失配。提出PF3plat，首个无需位姿的单次前馈3DGS框架。核心创新是解决像素对齐3DGS在无位姿/稀疏视角下3D高斯中心定位不准确导致的训练不稳定问题：通过预训练模型实现粗对齐，轻量级可学习模块精修，几何置信度评分评估3D高斯中心可靠性并条件化参数预测。在大规模真实数据集上全面SOTA。适用时还应保留论文给出的边界：几何置信度依赖预训练模型质量；极端宽基线（>90度）效果下降；动态场景不支持；前馈泛化受限于训练数据多样性。
% METHOD: 2410.24207
\paragraph{NoPoSplat（规范空间 pose-free 前馈）（2024.10）}
\textbf{锚点细化：} NoPoSplat：针对像素对齐方法``局部预测再变换''范式依赖精确位姿、变换-融合策略难以整合为连贯全局表示的局限，用``规范高斯空间直接预测''替代局部到全局的坐标变换，核心思想是将第一个输入视角的相机坐标系作为全局参考帧，所有视角的高斯基元直接预测到这一统一坐标系中、完全绕过坐标变换，纯 ViT 架构不注入任何几何先验，在低重叠场景下反而优于注入对极约束或代价体的专用骨干，且规范空间的高斯一致性反向解锁了位姿估计能力。
NoPoSplat~\cite{2410_24207} 是 pose-free 前馈 3DGS 的重要锚点。pixelSplat、MVSplat 和 DepthSplat 通常假设输入相机位姿已知：网络可以先在每个视图局部坐标里预测深度或高斯，再用外参把局部高斯变换到世界坐标。这个 ``预测再变换再融合'' 的流程在实际稀疏图像里很脆弱，因为位姿往往来自 SfM，低纹理、宽基线或少重叠都会让位姿误差和重建误差互相放大。NoPoSplat 的核心改变是干脆不输入外参，把第一张图像的局部相机坐标系固定为 canonical space，让所有视角的高斯都直接在这个规范空间里被预测出来。

模型结构故意保持简单：每张输入图像被切成 ViT token，相机内参被编码成一个 token 与图像 token 拼接，用来解决焦距引起的尺度模糊；编码器逐视图提取特征，解码器通过跨视角 attention 交换信息。输出端有两个 DPT 风格预测头：一个主要预测高斯中心，另一个在解码器特征之外再接入 RGB shortcut，预测尺度、旋转、不透明度和球谐颜色。训练时不需要真实深度或显式匹配损失，而是用目标视角光度渲染损失逼迫规范空间高斯同时解释多个输入视角。这样，跨视角对齐不再通过外部位姿矩阵完成，而是变成网络必须在 canonical space 中学会的融合问题。

NoPoSplat 还把重建结果反过来用于位姿估计：先从预测高斯中心得到粗位姿，再通过渲染图与输入图的光度误差做细化。相对 Splatt3R/DUSt3R 路线，它不把点云或深度监督作为必要条件；相对 PF3plat、FreeSplatter 等后续方法，它提供了一个更干净的参照点，即 ``规范空间高斯 + 内参 token + 光度训练''。这一设计解释了为什么它在输入图像重叠很小时仍能工作：模型不强行套用对极几何，而是先学习把多视图证据放进同一个局部参考系。\textbf{方法信息：} Google / DeepMind、NVIDIA、Microsoft 参与了这篇论文；代码：\url{https://github.com/cvg/NoPoSplat}；仓库提供安装、预训练权重、数据准备、训练和评测命令；GitHub 约 971 stars，Semantic Scholar 约 228 次引用。



% METHOD: 2411.17190
\paragraph{SelfSplat（无位姿无 3D 先验泛化 3DGS）（2024.11）}
SelfSplat~\cite{2411_17190} 不依赖已知位姿或 3D 先验。它用多视图和单目编码器提取特征，Swin Transformer 做跨视图交换，再由匹配感知位姿网络估计相对位姿。

位姿分支不是独立回归相机，而是利用匹配特征估计两帧间 SE(3) 变换；深度精化模块再用多视图几何一致性修正初始深度。高斯解码器根据精化后的几何输出像素对齐的位置、尺度、旋转、不透明度和颜色，目标视图只在训练中提供渲染监督。重投影损失和渲染损失共同更新位姿网络与高斯网络，使相机估计和重建质量互相促进。相对 NoPoSplat 的规范空间，SelfSplat 更强调“先估位姿、再用重建反过来改善位姿”的闭环。
\textbf{设计动机与边界：} 相对原始 3DGS 以独立高斯和图像重建误差为主的设定，SelfSplat 把问题落在以下可实现的改动上：技术路线：(1) 多视图和单目编码器：ResNet提取4x下采样特征，Swin Transformer进行跨视图信息交换。(2) 匹配感知位姿估计：专门设计的位姿网络估计两帧间SE(3)变换，通过匹配特征增强位姿精度。(3) 深度精化模块：利用多视图几何一致性精化初始深度估计，增强跨视图几何一致性。(4) 高斯解码器：预测像素对齐的高斯参数（位置、缩放、旋转、不透明度、颜色），仅从两个输入视图生成高斯，目标视图用于监督。(5) 联合训练：重投影损失和渲染损失联合训练位姿网络和高斯网络，位姿精度和重建质量互相促进。其设计目的不是单独增加一个模块，而是用上述表示、监督或渲染约束消除原流程在目标任务中的具体失配。首个无需位姿和3D先验的可泛化3DGS方法，通过匹配感知位姿估计和深度精化实现位姿精度与重建质量的互相促进。在RealEstate10K、ACID和DL3DV上取得SOTA，展现强跨数据集泛化能力。适用时还应保留论文给出的边界：假设已知相机内参；大基线多视图可能失败；位姿精度仍低于SfM方法
\textbf{方法信息：} 代码：\url{https://github.com/gynjn/selfsplat}；GitHub 关注度较高，约 147 stars。



% METHOD: 2412.09573
\paragraph{FreeSplatter（无标定稀疏视角同时重建与估计相机）（2024.12）}
\textbf{锚点细化：} FreeSplatter：针对 NoPoSplat 坐标规范化需精心设计训练损失的局限，用``流线型 Transformer + 两阶段退火训练''替代复杂训练策略，核心思想是预训练阶段同时使用深度位置监督损失和像素对齐余弦相似度损失提供稳定三维几何初始化，主训练阶段逐步退火位置监督并过渡到纯渲染损失，不仅重建质量超越依赖位姿的大规模重建模型，还能利用 PnP-RANSAC 从预测的高斯位置图高效估计相机参数。
FreeSplatter~\cite{2412_09573} 把无位姿进一步放宽到无标定输入：内参、外参都不给，模型仍需生成 3D Gaussians，并从结果中恢复相机参数。它的基本选择是把第一张输入图像作为参考相机坐标系，所有视角的高斯都预测到这个统一参考系中，从而避免无全局坐标时的任意性。

具体做法是把 $N$ 张图像切成 ViT 风格 token，加入位置嵌入和视图嵌入后，用多层自注意力交换跨视角信息，再解码为像素对齐 Gaussian map。每个像素输出一个 3D 高斯，参数包括位置、旋转、尺度、不透明度和球谐颜色。相机参数不是单独黑盒回归，而是从预测出的高斯位置图中用 PnP、Weiszfeld 等现成求解器恢复；这样相机估计与高斯几何共享同一个三维输出。论文同时给出物体级 FreeSplatter-O 和场景级 FreeSplatter-S，分别面向 Objaverse 物体和 RealEstate10K/MegaDepth 场景。

相对 NoPoSplat 的规范空间光度训练，FreeSplatter 更明确地把 ``高斯图'' 当作相机求解的中间几何对象；相对 GS-LRM，它去掉了对已知相机的依赖。由于每个像素的 2D 位置和预测 3D 位置天然成对，PnP 可以从 Gaussian map 中恢复相机；Weiszfeld 等求解步骤则帮助在多视角几何中获得更稳的参考关系。这个设计使重建和位姿估计共享同一显式 3D 输出，而不是训练两个彼此松散的头。\textbf{方法信息：} Tencent 参与了这篇论文；代码：\url{https://github.com/tencentarc/freesplatter}；仓库 README 提供环境、Gradio demo 和模型说明；GitHub 约 241 stars。



% METHOD: 2412.16604
\paragraph{OmniSplat（全景图像前馈高斯）（2024.12）}
OmniSplat~\cite{2412_16604} 处理等距柱状全景图输入。普通透视前馈模型在全景极区采样不均，特征提取和高斯预测都会产生伪影；OmniSplat 用阴阳网格把全景分解为两张准均匀图像。

预训练透视编码器在 Yin/Yang 域提取特征，阴阳交叉注意力对齐两个域；像素对齐高斯分别预测后，经阴阳 rasterization 在像素空间合并。相对 Pano2Room 的单全景室内重建，OmniSplat 更像把已有前馈模型适配到全景输入。
Yin 网格来自球坐标变换，Yang 网格通过旋转覆盖剩余区域，二者都比等距柱状投影更接近透视图像的均匀采样。预训练透视编码器因此可以继续使用，交叉注意力负责在两个域之间交换边界和重叠区域信息。每个域内的像素预测位置、协方差、不透明度和球谐颜色，再分别光栅化并在全景像素空间合并。关键增量是把“全景投影问题”转化为“两个准透视域的前馈高斯预测问题”。
\textbf{实现细节：} YoNoSplat（2025.11） 支持任意数量图像的输入，无论是否提供位姿和标定——模型预测每个视角的局部高斯和相机位姿，通过预测或提供的位姿聚合为全局表示。联合学习高斯和相机参数的核心困难在于两个任务的梯度纠缠。YoNoSplat 引入混合训练策略（mixing training strategy）缓解这一问题：初始阶段使用真值位姿聚合局部高斯（提供稳定的几何监督），逐步过渡到预测位姿与真值位姿的混合（防止训练不稳定和曝光偏差）。同时提出成对相机距离归一化（pairwise camera-distance normalization）解决尺度歧义——将相邻视角间的平移距离归一化，消除绝对尺度的不确定性；内参嵌入使模型能处理未标定输入（预测相机内参）。YoNoSplat 在 NVIDIA GH200 GPU 上仅需 2.69 秒即可从 100 张 $280 \times 518$ 分辨率图像完成场景重建。
\textbf{方法信息：} 代码：\url{https://github.com/esw0116/omnisplat}。
\textbf{设计动机与边界：} 相对原始 3DGS 以独立高斯和图像重建误差为主的设定，OmniSplat 把问题落在以下可实现的改动上：OmniSplat的核心方法：(1)阴阳网格分解：将全景等距柱状投影图像通过阴阳坐标系分解为两张准均匀图像。Yin网格定义为球坐标变换，Yang网格为剩余部分通过旋转矩阵获得。分解后的图像具有类似透视图像的均匀采样特性；(2)阴阳交叉注意力：在两个域之间进行特征对齐，准均匀特性使预训练透视编码器能正确提取特征；(3)像素对齐高斯预测：在每个像素位置预测3D高斯参数（位置、协方差、不透明度、球谐系数）；(4)阴阳光栅化：分别光栅化阴阳图像后在像素空间合并，避免全景渲染中的极地条纹伪影。整个过程无需额外训练，直接利用预训练的MVSplat或pixelSplat等前馈网络。其设计目的不是单独增加一个模块，而是用上述表示、监督或渲染约束消除原流程在目标任务中的具体失配。提出OmniSplat，首个面向全景图像的前馈3DGS生成方法。通过阴阳网格分解将全景图像转换为类透视的准均匀图像，利用预训练透视模型先验，无需训练即可实现高质量全景3DGS生成。适用时还应保留论文给出的边界：依赖预训练前馈网络的质量；阴阳网格分解在极地区域仍有轻微变形；对全景图像数量要求（两帧输入）。
% METHOD: 2412.16932
\paragraph{GSemSplat（无标定图像对语义高斯）（2024.12）}
GSemSplat~\cite{2412_16932} 建在 Splatt3R 之上，把 pose-free 3DGS 扩展到语义场。它新增语义头，为每个像素高斯输出区域特定特征和上下文感知特征。

底层仍是 Splatt3R 式权重共享 ViT 编码器、双 Transformer 解码器、高斯头和点云头，因此两张未标定图像即可预测像素级高斯。语义头绑定在这些高斯上：区域特定特征通过 SAM mask pooling 获得细粒度对象语义，上下文特征通过全图编码捕获低纹理区域的全局语义。两类特征再由 MLP 解压到 CLIP 空间，并用 2D CLIP 特征监督。相对 Uni3R 的多视图开放词汇，GSemSplat 更聚焦无标定图像对，把 Splatt3R 的几何前馈能力扩展成语义前馈能力。
\textbf{设计动机与边界：} 相对原始 3DGS 以独立高斯和图像重建误差为主的设定，GSemSplat 把问题落在以下可实现的改动上：GSemSplat的方法包括：(1)基干网络：基于Splatt3R架构（继承DUSt3R的pose-free特性），由权重共享ViT编码器、两个Transformer解码器、高斯头和点云头组成。给定两张未标定输入图像，预测每个像素的3D高斯参数；(2)语义头：新增语义头输出低维区域特定特征和上下文感知特征，绑定到对应像素预测的3D高斯上；(3)双特征方法：区域特定特征通过SAM mask pooling获得细粒度语义；上下文感知特征通过全图编码获得全局语义信息。两种特征互补，前者适合精确物体识别，后者适合低纹理区域；(4)训练与推理：语义特征通过MLP子网解压为高维CLIP特征，用2D CLIP特征监督训练。推理时单次前馈即可输出语义3D高斯。其设计目的不是单独增加一个模块，而是用上述表示、监督或渲染约束消除原流程在目标任务中的具体失配。提出GSemSplat，首个从稀疏未标定图像对进行泛化3D语义场建模的方法。构建在Splatt3R之上，通过双特征方法（区域特定+上下文感知语义特征）实现有效的3D语义学习，无需逐场景优化、密集图像或相机标定。适用时还应保留论文给出的边界：仅支持两视图输入；语义精度受限于CLIP特征的质量；对遮挡严重场景的处理有限。
\textbf{方法信息：} Microsoft 参与了这篇论文；代码：\url{https://github.com/wxrui182/gsemsplat}。



% METHOD: 2503.10286
\paragraph{VicaSplat（无姿态视频帧一次前向）（2025.03）}
VicaSplat~\cite{2503_10286} 面向无姿态视频帧。它用可学习 camera tokens 与视觉 tokens 双向注意力交互，同时预测 3D 高斯和相机姿态。

相机姿态用单位双四元数紧凑表示，并设计姿态对齐损失。训练从 DUSt3R/MASt3R 蒸馏空间知识，先学几何，再联合优化 NVS 和位姿回归。相对 SPFSplat 系列，VicaSplat 更强调视频帧序列和 camera token。
图像编码器把每帧变成视觉 tokens，camera tokens 与视觉 tokens 通过双向注意力交互：视觉 tokens 给相机分支提供匹配和外观线索，camera tokens 又把视角相关信息反馈给高斯预测头。不同预测头分别输出 3D 高斯属性和相机姿态参数。渐进训练先让模型从 DUSt3R/MASt3R 蒸馏到稳定空间几何，再联合训练新视角合成和姿态回归，避免一开始就让颜色重建误差牵动不稳定相机。
\textbf{实现细节：} NoPoSplat（2024.10） 对此提出了根本性的范式转变——不再"局部预测再变换"，而是让所有视角直接在同一个规范高斯空间（canonical Gaussian space）中预测高斯。具体做法是将第一个输入视角的相机坐标系作为全局参考帧，所有视角的高斯基元直接预测到这一统一坐标系中，完全绕过了局部到全局的坐标变换。形式化地，网络学习映射 $f_\theta: \{(I^v, K^v)\}_{v=1}^V \rightarrow \{(\mu_j^{v \rightarrow 1}, \alpha_j, r_j, s_j, c_j)\}$，其中上标 $v \rightarrow 1$ 表示第 $v$ 个视角中像素 $p_j$ 对应的高斯参数直接表达在第 1 帧的相机坐标系下，而非第 $v$ 帧的局部坐标系。第一个视角的位姿固定为 $[I \mid 0]$（单位旋转、零平移）。直觉上，这就像让所有画家面对同一幅画布各自画自己看到的部分，而不是每个人在自己的画布上画完再拼到一起。在架构层面，NoPoSplat 采用纯 Vision Transformer（ViT）编码器-解码器，不注入任何几何先验（如对极约束或代价体）。各视角图像经 patch 嵌入后与内参 token 拼接，分别送入共享权重的 ViT 编码器；解码器通过跨注意力层实现多视角信息交互。参数预测采用基于 DPT 架构的双头设计：(1) 位置头仅利用 Transformer 解码器特征预测高斯中心位置；(2) 属性头在解码器特征之外额外引入 RGB 图像作为输入（RGB shortcut），确保纹理细节信息直接流向高斯基元，弥补 ViT 特征因 16 倍下采样而丢失的高频结构信息。
\textbf{方法信息：} 代码：\url{https://github.com/wu-cvgl/vicasplat}。
\textbf{设计动机与边界：} 相对原始 3DGS 以独立高斯和图像重建误差为主的设定，VicaSplat 把问题落在以下可实现的改动上：VicaSplat的核心是定制的Transformer解码器：(1) 图像编码器将每帧映射为视觉令牌；(2) 可学习相机令牌与视觉令牌通过双向注意力交互，因果聚合多视角特征并注入视角相关特征；(3) 通过不同预测头生成3D高斯属性和相机姿态参数。使用单位双四元数紧凑表示相机姿态，开发新的姿态对齐损失函数。从DUSt3R/MASt3R预训练模型蒸馏空间知识而非从头训练。采用渐进训练策略：先聚焦场景几何估计，再联合优化NVS和相机姿态回归。其设计目的不是单独增加一个模块，而是用上述表示、监督或渲染约束消除原流程在目标任务中的具体失配。首个从无位姿视频帧一次前向传播同时预测3D高斯泼溅和相机姿态的框架，使用可学习相机令牌和双四元数表示，在跨数据集泛化上展现卓越能力。适用时还应保留论文给出的边界：对视频帧数量和时间连续性有要求；相机令牌的容量可能限制大规模场景；依赖预训练点云模型的知识蒸馏。
% METHOD: 2505.23716
\paragraph{AnySplat（无约束图像集合前馈 3DGS）（2025.05）}
AnySplat~\cite{2505_23716} 接在 NoPoSplat 之后处理更宽的输入假设：图像集合可以未校准，视角数量和覆盖范围也更自由。NoPoSplat 固定第一视角为规范空间，并假设内参可作为输入；AnySplat 则把相机估计和高斯预测一起放进一个 VGGT 风格几何 Transformer，目标是让模型从无约束图像集合中同时恢复相机、深度和 3D Gaussians。

它的主干先用 DINOv2 把每张图像切成 patch token，再加入 camera tokens 和 register tokens。Transformer 交替执行帧内 attention 和全局 attention：帧内 attention 保留单图局部结构，全局 attention 让不同视角交换几何证据。精炼后的 camera token 通过轻量 MLP 输出每个输入视角的内参和外参；图像 token 则进入高斯头和深度头，预测像素级高斯参数、深度图和相机相关几何。由于像素级高斯会在多视角重叠区域产生大量重复 primitive，AnySplat 进一步加入 differentiable voxelization，把相邻或同体素的像素高斯合并成 voxel-level Gaussians，在保持渲染质量的同时减少约 30--70\% 的高斯数量。

训练上，AnySplat 不要求 SfM/MVS 预处理出的真实相机，而是从预训练 VGGT 蒸馏相机和几何伪标签，再用渲染图像和深度监督约束最终高斯。相对 NoPoSplat，它更像把 ``无位姿重建'' 升级成 ``相机估计、深度估计、体素去冗余和高斯渲染'' 的联合系统；相对后续 ZipSplat，它仍以像素高斯为第一步，只是在体素化阶段压缩冗余。\textbf{方法信息：} 代码：\url{https://github.com/internrobotics/anysplat}；仓库公开训练和推理代码、Hugging Face 模型、demo 和评测脚本；GitHub 约 886 stars。
\textbf{设计动机与边界：} 相对原始 3DGS 以独立高斯和图像重建误差为主的设定，AnySplat 把问题落在以下可实现的改动上：AnySplat的架构包含：1)几何Transformer：遵循VGGT设计，使用DINOv2将每个图像分块为token，添加可学习相机token和寄存器token，通过交替注意力transformer（帧内注意力+全局注意力）处理所有视角的token；2)相机位姿预测：精修的相机token通过轻量MLP预测每个视角的相机参数，包括内参和外参；3)高斯参数预测：从几何Transformer输出的特征解码为高斯基元参数（位置、不透明度、协方差、颜色SH系数）；4)可微分体素化模块：将像素级高斯聚类到体素中，合并冗余基元，减少30-70%的高斯数量同时保持渲染质量，显著提升效率；5)伪标签知识蒸馏：从预训练VGGT骨干蒸馏相机位姿和几何先验，通过伪标签监督训练，无需3D标注或SfM/MVS预处理。其设计目的不是单独增加一个模块，而是用上述表示、监督或渲染约束消除原流程在目标任务中的具体失配。提出AnySplat，面向未校准图像集合的前馈新视角合成网络。单次前馈传递同时输出3D高斯基元（编码场景几何和外观）和每个输入图像的相机内参外参。引入可微分体素化模块，将像素级高斯合并为体素级高斯，消除30-70%冗余基元。设计伪标签知识蒸馏流程，从预训练VGGT骨干蒸馏相机和几何先验作为外部监督，无需任何3D SfM/MVS监督，仅依赖未校准图像训练。在稀疏和密集视角场景中均匹配位姿感知基线的质量。适用时还应保留论文给出的边界：前馈方法的精度上限受训练数据和模型容量限制；对极大规模场景的泛化能力有待验证；体素化可能损失精细细节。
% METHOD: 2506.09885
\paragraph{ULVS（更少 3D 依赖的大数据前馈 NVS）（2025.06）}
The Less You Depend, The More You Learn~\cite{2506_09885} 提出的 ULVS 不完全是一个 3DGS 表示方法，但它对 pose-free 前馈 3DGS 很重要，因为它挑战了本节许多方法的共同假设：显式 3D 表示、SfM 位姿和手工几何约束在小数据上有帮助，但在大规模数据上可能变成上限。论文把 pixelSplat、MVSplat、NoPoSplat 等偏几何约束的路线与 LVSM、RayZer 等数据驱动路线放在同一视角下比较，观察到依赖显式 3D 知识越少的方法，在训练数据增加时提升越快。

基于这个判断，ULVS 去掉两类依赖：不要求预定义的场景结构，也不要求训练或推理时的相机位姿标注。输入图像先由 DINOv2 tokenizer 转为 patch tokens，再由 Transformer 编码器和解码器建模场景 latent。难点是目标视角仍要可控，如果直接让网络用高维 latent 表示相机，latent 很容易泄漏目标图像内容。为此，Latent Learner 先把图像压缩成 7D 潜在位姿 token，即平移和四元数形式的低维瓶颈；再把这个紧凑 token 解析成像素级 ray embedding，作为解码器渲染条件。低维瓶颈限制信息泄漏，ray embedding 又保证每个像素仍有足够的视角条件。

因此，ULVS 相对 NoPoSplat 的差异不是把高斯放到另一个坐标系，而是把 ``显式 3DGS/位姿/SfM'' 的依赖整体降到最低，用大数据训练隐式 3D 感知。它给本章后续方法提供了另一种评价标准：当数据规模继续扩大时，是继续加入几何先验，还是把几何交给网络自己学。\textbf{方法信息：} Alibaba、ByteDance 参与了这篇论文；代码：\url{https://github.com/ou524u/Less3Depend}；仓库提供环境、数据处理、预训练模型下载、评测和训练命令；GitHub 约 66 stars。
\textbf{设计动机与边界：} 相对原始 3DGS 以独立高斯和图像重建误差为主的设定，The Less You Depend, The More You Learn 把问题落在以下可实现的改动上：ULVS基于对前馈NVS方法中偏向驱动(bias-driven)和数据驱动(data-centric)两种哲学的系统分析。发现减少显式3D知识依赖的方法在大规模数据下性能提升更快。ULVS框架：(1)基于Transformer的编码器-解码器架构，使用DINOv2编码输入图像为patch token；(2)Latent Learner核心组件，使用自编码器架构学习有意义的潜在位姿空间。编码器将图像蒸馏为高度紧凑的7D潜在位姿token（平移x和四元数q），低维瓶颈约束防止信息泄漏；解码器确保紧凑表示具有足够的渲染表达能力；(3)完全消除对场景结构和位姿标注的依赖，训练和推理均无需位姿信息，直接从大量2D图像学习隐式3D感知。其设计目的不是单独增加一个模块，而是用上述表示、监督或渲染约束消除原流程在目标任务中的具体失配。发现并验证了'依赖越少，学习越多'的规律——对3D知识依赖较少的方法在数据增加时性能加速更快，最终超越依赖3D知识的对手。据此设计无需场景结构和位姿标注的前馈NVS框架，达到SOTA甚至超越使用位姿训练数据的方法。适用时还应保留论文给出的边界：需要大量训练数据才能发挥优势；潜在位姿空间的可解释性有限；对极度稀疏输入（1张图像）的效果未明确。
% METHOD: 2507.22342
\paragraph{UFV-Splatter（不利视角适配）（2025.07）}
UFV-Splatter~\cite{2507_22342} 针对预训练 pose-free 模型在不利视角下失效的问题。所谓不利视角通常指目标不居中、尺度不合适或相机偏离训练分布。

方法先对图像重居中和平移缩放，把输入变成更接近有利视角；再用 LoRA 微调保留预训练先验；高斯适配器改善重居中图像生成高斯的几何一致性；高斯对齐方法对齐不同视角预测。相对 AnySplat 的通用无约束输入，UFV-Splatter 更像域适配工具。
训练上，它不需要重新收集大量不利视角数据，而是从有利视角数据集出发，用预训练模型合成偏心、偏尺度或偏姿态样本，再学习如何把这些输入拉回模型熟悉的分布。LoRA 只调整少量参数，避免破坏原前馈模型的几何先验；高斯适配器和对齐模块再修正重居中带来的跨视角几何偏差。相对 AnySplat 的通用无约束输入，UFV-Splatter 更像给已有模型加一个面向不利拍摄条件的适配层。
\textbf{实现细节：} FreeSplatter（2024.12） 采用流线型 Transformer 架构，自注意力块促进多视角图像 token 之间的信息交换，在统一参考帧中解码像素对齐的三维高斯基元。在训练策略上，FreeSplatter 设计了两阶段方案：预训练阶段同时使用深度位置监督损失（将预测高斯位置与深度解投影真值对齐）和像素对齐余弦相似度损失（约束高斯分布在相机射线上），为高斯位置学习提供稳定的三维几何初始化；主训练阶段逐步退火位置监督并过渡到纯渲染损失加像素对齐约束，使模型最终以渲染质量为目标优化。此外，FreeSplatter 简化了文本/图像到 3D 的内容创建管线，消除了相机位姿管理的复杂性。
\textbf{方法信息：} 代码：\url{https://github.com/yfujimura/ufv-splatter}。
\textbf{设计动机与边界：} 相对原始 3DGS 以独立高斯和图像重建误差为主的设定，UFV-Splatter 把问题落在以下可实现的改动上：UFV-Splatter解决预训练模型在不利视角（相机未对准物体中心）下表现差的问题。方法包括：(1)图像重居中：通过简单图像平移和缩放将不利视角图像转换为有利视角格式；(2)LoRA微调：在预训练模型中添加低秩适应层保持原有先验；(3)高斯适配器模块：改善重居中图像生成高斯的几何一致性；(4)高斯对齐方法：对齐不同视角的预测以实现准确的目标视图渲染；(5)训练策略：仅使用有利视角数据集（如G-Objaverse），利用预训练模型合成不利视角训练图像。其设计目的不是单独增加一个模块，而是用上述表示、监督或渲染约束消除原流程在目标任务中的具体失配。提出首个使预训练无姿态前馈3DGS模型适应不利视角的框架。通过图像重居中+LoRA微调利用有利视角的先验，设计高斯适配器模块增强几何一致性，提出高斯对齐方法和仅需有利图像的训练策略。适用时还应保留论文给出的边界：对极端视角偏移的效果可能下降；依赖预训练模型的质量；主要针对物体中心场景。
% METHOD: 2508.03643
\paragraph{Uni3R（重建与开放词汇语义统一）（2025.08）}
Uni3R~\cite{2508_03643} 在 pose-free 前馈重建中加入开放词汇语义理解。它用 DINOv2 提取 patch 特征，跨视角 Transformer 融合多视图信息，DPT 解码器和 MLP 头同时预测高斯参数、深度和语义特征。

相机内参嵌入编码焦距和主点，用来缓解尺度歧义；跨视角 Transformer 同时做帧内自注意力和帧间全局注意力，使任意数量输入视角都能交换几何证据并保持排列等变。点图引导几何损失来自冻结 VGGT，带置信度的密集点图作为软监督，防止前馈高斯陷入局部几何最优。语义特征则随高斯一起输出，可被文本查询用于 3D open-vocabulary segmentation。相对 GSemSplat 的图像对语义，Uni3R 支持多视图，并把新视角合成、深度预测和开放词汇理解放进同一次前向。
\textbf{设计动机与边界：} 相对原始 3DGS 以独立高斯和图像重建误差为主的设定，Uni3R 把问题落在以下可实现的改动上：Uni3R采用以下架构：(1)内参嵌入：编码相机焦距和主点解决单目重建的尺度模糊；(2)跨视角Transformer编码器：基于DINOv2提取每视角patch特征，通过帧内自注意力和帧间全局注意力融合多视角信息，支持任意多视角输入并保持排列等变性；(3)DPT解码器+MLP头：预测像素对齐的3D高斯参数，包括语义特征向量；(4)点图引导几何损失：使用冻结的VGGT生成带置信度的密集点图作为软监督，防止模型陷入局部最优。语义特征支持零样本文本查询实现3D语义分割。其设计目的不是单独增加一个模块，而是用上述表示、监督或渲染约束消除原流程在目标任务中的具体失配。提出首个从无位姿多视角图像同时实现3D重建和开放词汇语义理解的前馈框架。跨视角Transformer实现全局一致的3D表示，单一前馈过程即可完成新视角合成、开放词汇3D语义分割和深度预测。在RE10K达到25.07 PSNR，ScanNet达到55.84 mIoU。适用时还应保留论文给出的边界：依赖DINOv2和VGGT等大型预训练模型，计算资源需求较高。语义特征的分辨率受限于patch级特征提取。
\textbf{方法信息：} 代码：\url{https://github.com/horizonrobotics/uni3r}；GitHub 关注度较高，约 187 stars。



% METHOD: 2511.07321
\paragraph{YoNoSplat（有无位姿统一模型）（2025.11）}
\textbf{锚点细化：} YoNoSplat：支持任意数量图像的输入，无论是否提供位姿和标定——模型预测每个视角的局部高斯和相机位姿，通过预测或提供的位姿聚合为全局表示。联合学习高斯和相机参数的核心困难在于两个任务的梯度纠缠。YoNoSplat 引入混合训练策略（mixing training strategy）缓解这一问题：初始阶段使用真值位姿聚合局部高斯（提供稳定的几何监督），逐步过渡到预测位姿与真值位姿的混合（防止训练不稳定和曝光偏差）。同时提出成对相机距离归一化（pairwise camera-distance normalization）解决尺度歧义——将相邻视角间的平移距离归一化，消除绝对尺度的不确定性；内参嵌入使模型能处理未标定输入（预测相机内参）。YoNoSplat 在 NVIDIA GH200 GPU 上仅需 2.69 秒即可从 100 张 $280 \times 518$ 分辨率图像完成场景重建。
YoNoSplat~\cite{2511_07321} 进一步把 pose-free 问题扩展成统一输入条件问题：同一个模型要能处理任意数量图像，也要能在有位姿、无位姿、有内参、无内参之间切换。它没有把所有图像强行映射到一个固定规范空间，而是采用局部到全局范式：每个输入视角先预测自己的局部高斯和相机参数，再用提供或预测的位姿把局部高斯聚合到全局坐标系。

难点在于高斯和相机参数会互相拖累。相机错会让高斯聚合错，高斯渲染错又会给相机分支错误梯度。YoNoSplat 的 mixing training strategy 先用真值位姿聚合局部高斯，让模型建立稳定几何；随后逐步混入预测位姿，让训练分布接近测试时真正会遇到的无位姿状态。尺度歧义通过成对相机距离归一化处理，即把相邻视角间的平移尺度标准化；内参嵌入网络则让模型在未标定输入下预测或吸收相机内参。

与 PF3plat 借助外部几何模型先粗对齐不同，YoNoSplat 的锚点意义是统一接口和可扩展视角数：一套前馈网络可以从少量图像到 100 张图像工作，也可以在给定位姿时直接使用真值位姿、无位姿时改用预测位姿。局部高斯先在各自相机系内预测，再通过位姿聚合，因此新增视角主要增加局部分支计算，而不是要求单一规范空间一次吸收所有视图；这也更接近真实照片集合的增量接入方式。论文报告在 NVIDIA GH200 上从 100 张 $280\times518$ 图像重建约 2.69 秒。\textbf{方法信息：} Microsoft 参与了这篇论文；代码：\url{https://github.com/cvg/YoNoSplat}；仓库提供预训练检查点、评测命令和训练配置；GitHub 约 210 stars。

\textbf{实现细节：} YoNoSplat（2025.11） 支持任意数量图像的输入，无论是否提供位姿和标定——模型预测每个视角的局部高斯和相机位姿，通过预测或提供的位姿聚合为全局表示。联合学习高斯和相机参数的核心困难在于两个任务的梯度纠缠。YoNoSplat 引入混合训练策略（mixing training strategy）缓解这一问题：初始阶段使用真值位姿聚合局部高斯（提供稳定的几何监督），逐步过渡到预测位姿与真值位姿的混合（防止训练不稳定和曝光偏差）。同时提出成对相机距离归一化（pairwise camera-distance normalization）解决尺度歧义——将相邻视角间的平移距离归一化，消除绝对尺度的不确定性；内参嵌入使模型能处理未标定输入（预测相机内参）。YoNoSplat 在 NVIDIA GH200 GPU 上仅需 2.69 秒即可从 100 张 $280 \times 518$ 分辨率图像完成场景重建。



% METHOD: 2604.21182
\paragraph{WildSplatter（无约束图像与外观控制）（2026.04）}
WildSplatter~\cite{2604_21182} 面向野外图像集：未知姿态、未知内参、光照变化都存在。它用 Depth Anything 3 ViT 骨干，一部分局部注意力提取帧特征，另一部分全局注意力跨帧聚合。

几何分支预测无色高斯，假设几何相对光照不变；外观分支从目标图像估计低维 appearance embedding，与局部场景特征一起预测颜色。这样同一几何可以在不同光照嵌入下渲染不同外观。相对 GSemSplat 的语义场，WildSplatter 关注外观控制。
\textbf{实现细节：} WildSplatter（2026.04） 处理无约束照片集合中两个根本性挑战：未知相机参数和变化光照条件。其核心设计是联合学习三维高斯和外观嵌入（appearance embeddings）——每个高斯除几何和颜色属性外，还携带一个以输入图像为条件的外观嵌入，支持在不同光照条件下灵活调制高斯颜色。具体而言，外观嵌入与空间位置联合条件化颜色预测：同一三维位置在不同光照条件下通过不同的外观嵌入产生不同颜色，而解耦的外观嵌入确保光照变化不影响几何表示的一致性。
\textbf{技术实现补充：} WildSplatter（2026.04） 处理无约束照片集合中两个根本性挑战：未知相机参数和变化光照条件。其核心设计是联合学习三维高斯和外观嵌入（appearance embeddings）——每个高斯除几何和颜色属性外，还携带一个以输入图像为条件的外观嵌入，支持在不同光照条件下灵活调制高斯颜色。具体而言，外观嵌入与空间位置联合条件化颜色预测：同一三维位置在不同光照条件下通过不同的外观嵌入产生不同颜色，而解耦的外观嵌入确保光照变化不影响几何表示的一致性。
\textbf{方法信息：} 代码：\url{https://github.com/yfujimura/wildsplatter}。
\textbf{设计动机与边界：} 相对原始 3DGS 以独立高斯和图像重建误差为主的设定，WildSplatter 把问题落在以下可实现的改动上：WildSplatter的架构基于Depth Anything 3 ViT编码器：(1)ViT骨干：分为两部分，第一部分用局部注意力提取帧级特征，第二部分用全局注意力跨帧聚合；(2)几何估计：从上下文图像预测无色3D高斯（中心、不透明度、旋转、尺度），基于光照不变性假设；(3)外观建模：从目标图像估计全局外观嵌入，与局部场景特征结合预测高斯颜色，外观嵌入在低维潜空间中捕获光照变化；(4)训练：在野外照片集（如Instagram收集的图像）上训练，使用光度损失和深度一致性损失。外观嵌入可跨场景泛化，允许在新场景中通过指定嵌入控制渲染光照。其设计目的不是单独增加一个模块，而是用上述表示、监督或渲染约束消除原流程在目标任务中的具体失配。提出首个面向无约束图像（未知姿态+变化光照）的前馈3DGS模型WildSplatter，通过几何-外观解耦和外观嵌入实现不到1秒的3D重建，同时支持多光照条件下的外观控制。适用时还应保留论文给出的边界：对极端光照变化（如夜景到日光）的外观控制可能有限；前馈模型的重建质量可能低于逐场景优化方法；依赖训练数据中光照变化的多样性。

\section{大规模重建模型}

LRM 路线把前馈 3DGS 放进更大规模的 transformer/Mamba/TTT 模型中。GS-LRM 是早期像素高斯 LRM，Long-LRM 解决长序列，Long-LRM++ 用半显式特征高斯保细节，C3G 用少量查询 token 建紧凑表示，tttLRM 则用 test-time training 层作为长上下文神经记忆。
% METHOD: 2404.19702
\paragraph{GS-LRM（像素高斯大重建模型）（2024.04）}
GS-LRM~\cite{2404_19702} 是 LRM 路线接入 3D Gaussian Splatting 的早期锚点。传统 LRM 多把少量输入图像编码成固定分辨率三平面，再由隐式场或解码器渲染目标视角；它的缺点是输出分辨率和三平面网格绑定，细节容易被固定 latent grid 抹平。原始 3DGS 则有显式、高频、实时渲染的优势，但需要逐场景优化。GS-LRM 的改变是把 LRM 的大规模前馈能力和 3DGS 的显式渲染表示直接接起来：输入少量带位姿图像，Transformer 一次输出像素对齐的高斯集合。

具体流程是，模型把 RGB 图像与 Plucker ray 坐标拼成 9 通道输入。Plucker ray 同时携带相机位姿和像素射线方向，因此 patchify 后不再需要额外的位置嵌入来告诉模型每个 token 对应哪条空间射线。多视图 patch tokens 进入一串 Transformer block，自注意力负责在不同视角之间交换几何和外观证据；输出 token 经线性头解码为每像素高斯参数，包括中心位置、尺度、旋转、不透明度和球谐颜色。因为一个输入像素可以直接产生一个或一组高斯，输出分辨率随输入图像变化，而不是被固定三平面分辨率限制。

GS-LRM 的训练也服务于这个目标：先在低分辨率上学习稳定几何，再提高分辨率补外观细节；Objaverse 物体和 RealEstate10K 场景分别训练，使同一框架能覆盖物体生成和真实场景重建。相对 pixelSplat/MVSplat，它更像大模型化的前馈重建；相对 Long-LRM、Long-LRM++，它是后续长序列和半显式特征高斯的基线参照。\textbf{方法信息：} Adobe 参与了这篇论文；引用很多，Semantic Scholar 约 340 次；本轮未发现官方完整代码。
\textbf{设计动机与边界：} 相对原始 3DGS 以独立高斯和图像重建误差为主的设定，GS-LRM 把问题落在以下可实现的改动上：GS-LRM的核心架构：(1)将输入的姿态图像通过Plucker射线坐标进行姿态编码（9通道特征图）；(2)分块化（patchify）为patch token（无额外位置嵌入，Plucker坐标自然提供空间信息）；(3)多视图token通过Transformer块序列（自注意力+MLP）处理；(4)线性层从输出token解码每像素高斯参数（位置、缩放、旋转、不透明度、SH颜色）。预测的每像素高斯沿光线放置，自然适应输入分辨率，保留高频细节。与之前LRM使用固定分辨率三平面不同，GS-LRM的输出分辨率随输入变化。训练采用分阶段策略，先低分辨率后高分辨率。其设计目的不是单独增加一个模块，而是用上述表示、监督或渲染约束消除原流程在目标任务中的具体失配。提出极简的Transformer架构从2-4张稀疏图像在0.23秒内预测高质量3D高斯；采用每像素高斯预测策略，自然处理不同尺度和复杂度的场景；在Objaverse和RealEstate10K上分别训练，均大幅超越SOTA基线；支持下游3D生成任务。适用时还应保留论文给出的边界：每像素高斯数量固定导致背景区域冗余；大场景需要更多输入视图；训练数据需求大。
% METHOD: 2410.12781
\paragraph{Long-LRM（长序列高斯重建模型）（2024.10）}
Long-LRM~\cite{2410_12781} 扩展 GS-LRM 到 32 张高分辨率输入。问题是多图像 patch token 可达数十万，纯 Transformer 成本过高；Long-LRM 使用 Mamba2 和少量 Transformer 的混合架构。

多视角图像和 Plucker 射线先被拼接并 patch 化，成为从序列到像素对齐高斯参数的映射问题。Mamba2 块以线性复杂度处理长序列，双向扫描让每个 token 同时吸收前后文；少量 Transformer 块补充全局自注意力，token 合并模块降低中间计算，训练和测试时的高斯剪枝控制输出大小。相对 GS-LRM，它的主要增量是长上下文和宽覆盖场景：把前馈高斯重建从少量视图推进到几十张高分辨率图像。
\textbf{设计动机与边界：} 相对原始 3DGS 以独立高斯和图像重建误差为主的设定，Long-LRM 把问题落在以下可实现的改动上：方法核心是处理32张960x540图像产生的约250K token长序列：（1）混合架构：7个Mamba2块+1个Transformer块的混合配置，Mamba2块以O(L)复杂度处理长序列，Transformer块提供全局自注意力。（2）双向扫描：对Mamba2采用双向扫描（前向+后向），更适合图像token处理。（3）Token合并模块：在混合块之间注入轻量token合并模块减少计算量。（4）高斯剪枝：训练和测试时的高斯剪枝提升高分辨率设置的效率。将多视角图像和Plucker射线拼接后patch化，视为序列到序列映射回归像素对齐高斯参数。其设计目的不是单独增加一个模块，而是用上述表示、监督或渲染约束消除原流程在目标任务中的具体失配。提出Long-LRM，首个支持32张输入图像的前馈3DGS重建模型，1秒内完成重建（比优化方法快800倍），输入量比前馈方法大60倍以上，质量可比优化方法。适用时还应保留论文给出的边界：需要A100 80G GPU；对更大规模场景（超过32张图）的扩展性待验证；训练数据需求大
\textbf{方法信息：} Adobe、Apple 参与了这篇论文；引用很多，Semantic Scholar 约 101 次。



% METHOD: 2512.04021
\paragraph{C3G（2K 查询 token 紧凑表示）（2025.12）}
C3G~\cite{2512_04021} 用约 2K 个 learnable query tokens 替代逐像素高斯预测。每个查询 token 通过自注意力与多视角 VGGT 特征交互，自动关注场景某一区域并解码成高斯。

注意力模式还可用于把 DINOv2 等 2D 语义特征提升到 3D，实现开放词汇分割和视角不变特征聚合。相对 GS-LRM 的像素高斯，C3G 证明较少、语义更明确的 token 也能支撑高质量重建。
具体来说，$N$ 个查询 token $Q\in\mathbb{R}^{N\times d}$ 与 VGGT 提取的多视角特征 $F$ 拼接成统一序列 $X=[Q;F]$，经过多层自注意力后，每个查询 token 学会从任意视角、任意区域读取相关证据。它不再把像素到高斯写成刚性一一映射，而是让 token 自行发现场景区域并聚合跨视角特征。这个查询机制同时让模型容易绑定语义：同一 token 关注的区域可被 DINOv2/CLIP 特征提升成 3D 语义表示。
\textbf{实现细节：} MVP 通过双层次注意力提升了可扩展性，但在高斯数量上仍采用密集像素对齐。C3G（2025.12） 从表示效率的角度出发，提出与 GS-LRM/Long-LRM 的密集像素对齐高斯根本不同的紧凑范式——仅用 $N{=}2048$ 个可学习查询 token 经 Transformer 自注意力聚合多视角特征后，每个 token 解码为一个三维高斯。与 §相应公式 节中 TokenGS 的可学习 token 方案类似，C3G 的核心优势是高斯数量与输入分辨率完全解耦，仅 2K 高斯即可表征完整场景。C3G 采用 VGGT 作为默认视觉编码器，训练中引入渐进低通滤波策略（RAIN-GS，从 $s{=}300$ 退火至 $s{=}0.3$）稳定初期位置预测。令人惊讶的是，仅通过光度重建损失训练，C3G 的查询 token 自发学会关注几何有意义的区域——注意力图显示每个 token 精准对齐到多视角中的同一物体区域，形成了无监督的跨视角对应关系。
\textbf{方法信息：} Sony 参与了这篇论文。
\textbf{设计动机与边界：} 相对原始 3DGS 以独立高斯和图像重建误差为主的设定，C3G 把问题落在以下可实现的改动上：C3G框架包含三个核心设计：（1）可学习查询Token机制：引入N个可学习查询Token Q∈R\^{}{N×d}，每个Token负责发现和表示3D场景的特定区域，通过自注意力与多视角图像特征交互，灵活关注任意视角的任意区域，而非像PixelSplat等方法刚性映射每个像素到高斯。（2）多视角注意力聚合：将查询Token与视觉编码器（VGGT）提取的多视角特征拼接为统一序列X=[Q;F]，通过多层自注意力Transformer处理，查询Token学习聚合相关跨视角视觉特征形成连贯3D高斯。（3）注意力引导的特征提升：利用查询Token学习到的注意力模式用于高斯解码，高效将2D特征（如DINOv2语义特征）提升到3D空间，实现场景理解。训练仅使用光度重建损失，无需深度或场景分解的显式监督，每个Token自然学会表示不同区域。该方法在保持高重建质量的同时，将高斯数量从数十万降至2K，显著降低内存开销。其设计目的不是单独增加一个模块，而是用上述表示、监督或渲染约束消除原流程在目标任务中的具体失配。提出C3G框架，用仅2K个可学习查询Token替代逐像素高斯预测，实现内存高效且几何意义明确的紧凑3D表示。通过自注意力聚合多视角特征指导高斯生成，在无位姿新视角合成、3D开放词汇分割和视角不变特征聚合三个任务上达到最优。证明了紧凑且几何有意义的表示足以实现高质量场景重建和理解。适用时还应保留论文给出的边界：仅使用2K个高斯可能在极其复杂的场景中丢失细节；前馈方式可能对训练数据分布外的场景泛化能力有限；依赖VGGT作为视觉编码器，继承了其计算开销。
% METHOD: 2512.10267
\paragraph{Long-LRM++（半显式特征高斯）（2025.12）}
Long-LRM++~\cite{2512_10267} 发现显式高斯在前馈 LRM 中容易细节模糊，而隐式方法又渲染慢。它用半显式表示：每个 token 预测 feature Gaussians，把特征 splat 到目标画布，再由轻量解码器输出 RGB。

骨干仍类似 Long-LRM 的 Mamba2/Transformer 混合结构，但因为只预测特征高斯，可用更少高斯。目标帧解码器把 3D splatting 的低频信息与输入视角高频特征融合。相对 Long-LRM，它不是直接渲染颜色，而是渲染可解码特征。
\textbf{实现细节：} 其后续工作 Long-LRM++（2025.12） 在 Long-LRM 的基础上将逐像素颜色高斯替换为半显式的"特征高斯"表示，在保持前馈效率的同时显著提升细节重建质量。核心改动在于：每个图像 token 预测 $K$ 个自由移动的特征高斯（而非射线对齐的颜色高斯），每个高斯携带 $F$ 维特征向量而非 SH 颜色系数。渲染时，特征高斯通过 splatting 将特征渲染到目标视角的特征画布上，再与目标相机的 Pl\"ucker 射线坐标拼接后送入轻量级 Transformer 解码器（平移不变局部注意力捕获高频细节 + 全局注意力建模长程依赖），最终线性层输出 RGB 颜色。这一设计的核心优势在于：(1) 特征高斯的自由移动特性提供比像素对齐更强的鲁棒性和表达力——无需精确对齐像素即可编码场景信息；(2) 解码器的隐式表达能力弥补了显式高斯在细粒度纹理上的不足。此外，Long-LRM++ 引入多空间分割合并机制：将高斯分为多组独立渲染解码，再通过自适应加权融合各组结果，有效处理镜面反射等非物理几何场景。超参数分析显示，每 token 预测的高斯数 $K$ 和特征维度 $F$ 对质量影响显著——增大 $K$ 或 $F$ 允许模型存储更多场景信息，但计算开销线性增加。迭代精炼与可扩展架构。
\textbf{方法信息：} Adobe 参与了这篇论文。
\textbf{设计动机与边界：} 相对原始 3DGS 以独立高斯和图像重建误差为主的设定，Long-LRM++ 把问题落在以下可实现的改动上：Long-LRM++解决显式GS方法的细节模糊和隐式方法的渲染速度限制：(1) 输入处理骨干：采用类似Long-LRM的Mamba2+Transformer混合骨干，1T7M3块序列，隐藏维度768，移除token合并模块。由于使用半显式表示，仅需1/4数量的高斯（相比Long-LRM），显著降低GPU内存。(2) 特征高斯预测和溅射：每个token预测K个高斯，每个高斯关联特征向量f∈F\^{}F、2D偏移o∈R\^{}2、射线对齐深度及其他属性（旋转、不透明度、缩放）。2D偏移允许高斯自由移动。将高斯特征溅射到画布（尺寸为输入分辨率的1/r），拼接目标相机的Plucker射线。(3) 目标帧解码器：渲染的特征图通过线性层和若干轻量Transformer块解码为最终RGB。解码器将3D溅射的低频信息与输入视角的高频信息融合。这避免了隐式方法每帧需要24层网络前向传播的开销。(4) 训练：端到端训练，支持深度监督。其设计目的不是单独增加一个模块，而是用上述表示、监督或渲染约束消除原流程在目标任务中的具体失配。提出Long-LRM++，采用半显式场景表示（特征高斯+轻量解码器），在DL3DV上匹配LaCT的渲染质量，同时实现14 FPS实时渲染。仅需1/4数量的高斯，支持64输入视角的950x540分辨率，在ScanNetv2上新视角深度预测也优于直接高斯深度渲染。适用时还应保留论文给出的边界：解码器仍需额外计算；对输入相机姿态误差敏感；更高分辨率的扩展性有待验证。
% METHOD: 2602.20160
\paragraph{tttLRM（Test-Time Training 长上下文记忆）（2026.02）}
tttLRM~\cite{2602_20160} 使用 LaCT/TTT 层作为长上下文神经记忆。多视角图像和光线嵌入被写入 TTT 层的快速权重中，形成隐式 3D 表示；虚拟 token 再查询这些权重，解码为 3DGS 等显式格式。

模型既支持批量前馈，也支持在线流式输入的渐进重建。它还从预训练隐式视图合成模型蒸馏知识，只微调解码器即可输出显式高斯。相对 Long-LRM 的序列模型，tttLRM 的锚点是“测试时训练层即记忆”。
LaCT 层把输入 token 的 key-value 对当作训练样本，使用 Muon 优化器在前向过程中更新快速权重 $W$，以固定大小神经记忆编码 KV 缓存，从而把长序列处理保持在线性复杂度。给定 $N$ 张带位姿图像，图像 token 与光线嵌入分块进入 LaCT；虚拟 token 再作为查询从快速权重中检索信息，在 3DGS 重建中对应虚拟视角或高斯属性预测入口。每个 LaCT 层还含窗口注意力，保留视角内局部结构。相对 Long-LRM 的静态长序列编码，tttLRM 把测试时适配本身变成模型记忆机制。
\textbf{实现细节：} tttLRM（2026.02） 从另一个维度突破输入规模限制——引入测试时训练（Test-Time Training, TTT）层，将多个图像观测高效压缩到 TTT 层的快速权重中，形成潜在空间中的隐式三维表示。其核心思想与 Long-LRM 的 Mamba2 状态空间模型有本质区别：Mamba2 将历史信息压缩到固定维度的隐状态中（递推式更新），而 tttLRM 的 LaCT（Large Chunk Test-Time Training）层将输入 token 的键-值对视为训练数据，用 Muon 优化器在线更新快速权重矩阵 $W$（梯度下降式更新），能够一次处理多达百万 token 的大块输入。前馈重建模式。模型由 24 个 LaCT 块组成（隐藏维度 768），输入图像与 Pl\"ucker 射线嵌入拼接后分块化为 token。LaCT 块包含窗口注意力（捕获视角内局部关系）和快速权重更新/查询操作（均为线性复杂度）。虚拟视角 token 查询快速权重获取更新后的表示，再经线性解码器预测逐像素高斯参数。训练损失为 MSE + VGG-19 感知损失，配合深度正则化和 opacity 正则化（与 Long-LRM 类似）。自回归重建模式。tttLRM 的独特能力在于支持流式输入——新视角以 mini-batch 方式逐批到达，快速权重增量更新，每步立即预测当前的三维高斯。这一模式模拟了机器人在环境中逐步探索并实时构建三维地图的场景。分布式序列并行策略允许多 GPU 分片处理长序列，进一步扩展输入规模。
\textbf{技术实现补充：} tttLRM（2026.02） 从另一个维度突破输入规模限制——引入测试时训练（Test-Time Training, TTT）层，将多个图像观测高效压缩到 TTT 层的快速权重中，形成潜在空间中的隐式三维表示。其核心思想与 Long-LRM 的 Mamba2 状态空间模型有本质区别：Mamba2 将历史信息压缩到固定维度的隐状态中（递推式更新），而 tttLRM 的 LaCT（Large Chunk Test-Time Training）层将输入 token 的键-值对视为训练数据，用 Muon 优化器在线更新快速权重矩阵 $W$（梯度下降式更新），能够一次处理多达百万 token 的大块输入。前馈重建模式。模型由 24 个 LaCT 块组成（隐藏维度 768），输入图像与 Pl\"ucker 射线嵌入拼接后分块化为 token。LaCT 块包含窗口注意力（捕获视角内局部关系）和快速权重更新/查询操作（均为线性复杂度）。虚拟视角 token 查询快速权重获取更新后的表示，再经线性解码器预测逐像素高斯参数。训练损失为 MSE + VGG-19 感知损失，配合深度正则化和 opacity 正则化（与 Long-LRM 类似）。自回归重建模式。tttLRM 的独特能力在于支持流式输入——新视角以 mini-batch 方式逐批到达，快速权重增量更新，每步立即预测当前的三维高斯。这一模式模拟了机器人在环境中逐步探索并实时构建三维地图的场景。分布式序列并行策略允许多 GPU 分片处理长序列，进一步扩展输入规模。
\textbf{方法信息：} Adobe 参与了这篇论文。
\textbf{设计动机与边界：} 相对原始 3DGS 以独立高斯和图像重建误差为主的设定，tttLRM 把问题落在以下可实现的改动上：tttLRM基于LaCT（Large Chunk Test-Time Training）层构建，该层将输入token的key-value对作为训练数据更新快速权重W，以固定大小的神经记忆编码KV缓存，实现线性复杂度。给定N张带位姿的图像，将其与光线嵌入拼接后分块并token化，通过LaCT块使用Muon优化器迭代更新快速权重。引入虚拟token作为查询来检索信息，在3DGS重建中这些是虚拟视角用于预测高斯属性。每个LaCT层包含窗口注意力模块捕捉视角内局部关系。模型支持两种模式：批量前馈模式一次性处理所有输入视角，以及在线学习模式支持流式观测的渐进重建。此外，利用预训练的隐式潜空间3D模型（如LaCT-based视图合成模型）的知识蒸馏，通过微调解码器层将隐式表示转换为显式3DGS输出，显著提升重建质量和收敛速度。其设计目的不是单独增加一个模块，而是用上述表示、监督或渲染约束消除原流程在目标任务中的具体失配。提出tttLRM，利用Test-Time Training (TTT)层实现线性复杂度的长上下文自回归3D重建。框架将多视角图像压缩到TTT层的快速权重中形成隐式3D表示，可解码为3DGS等显式格式。在线学习变体支持流式输入的渐进式重建。证明在新视角合成任务上的预训练可有效迁移到显式3D建模，在物体和场景的3DGS重建上取得SOTA性能。适用时还应保留论文给出的边界：主要聚焦前馈重建，对极端稀疏视角（如1-2视角）的场景仍有提升空间；在线学习模式需要足够的观测积累才能达到高质量；解码器从隐式到显式的转换可能在复杂几何细节上存在损失。

\section{单图到 3D 前馈}

单图前馈 3DGS 的困难是缺少多视角约束，尺度和遮挡背面都不确定。Gamba、studentSplat 和 CATSplat 分别用 Mamba、高质量教师和文本/空间双先验来补这一缺口。
% METHOD: 2403.18795
\paragraph{Gamba（Mamba 单图 3DGS）（2024.03）}
\textbf{锚点细化：} Gamba：针对此前 LRM 系列依赖 Transformer 自注意力（复杂度随序列长度二次增长）、需两阶段训练（先预测 4096 高斯再超分辨至 16384）的局限，用``Mamba 状态空间模型线性序列建模''替代 Transformer 自注意力，核心思想是利用 Mamba 的 $O(L)$ 线性复杂度直接预测 $L{=}16384$ 个三维高斯——状态空间模型（SSM）将整段历史信息压缩到一个固定维度的隐状态向量中，新 token 只需与当前隐状态交互（递推式 $h_t = \mathbf{A}h_{t-1} + \mathbf{B}x_t$），因此计算量与序列长度成正比而非二次增长，配合径向掩码约束（对投影中心落在物体轮廓之外的高斯施加约束）替代三维点云监督、防止模型坍塌，在 0.05 秒内完成重建、比 SDS 优化方法快约 1000 倍。
Gamba~\cite{2403_18795} 是单图物体级 3DGS 的早期锚点。单图重建没有多视角约束，传统 SDS 优化方法可以通过反复优化补信息，但耗时长；LRM 类前馈方法速度快，却常用 Transformer 处理大量 3D token，序列长度一长计算就变重。Gamba 的核心改变是把单图到高斯的预测写成序列建模问题，并用 Mamba 状态空间模型替代二次复杂度自注意力。

方法先用 DINO 编码参考图像，得到条件 token；再把这些 token 与一组可学习 3DGS token 拼接，输入 GambaFormer。Mamba 层通过状态空间递推把历史信息压缩到隐状态中，序列长度增长时计算近似线性，因此可以直接处理大量高斯 token。输出端为每个 token 解码高斯位置、尺度、旋转、不透明度和颜色，并通过 3DGS 可微 rasterizer 与多视角监督图像计算重建损失。

单图高斯训练还容易出现空间坍塌：网络可能把高斯放到不合理位置但仍在某些视角上凑出颜色。Gamba 用从多视角 mask 推导的 radial mask constraints 约束高斯投影分布，减少对三维点云预热监督的依赖；训练上采用渐进策略，先让较少或较低分辨率的高斯稳定学习，再逐步增加表示容量。相对 GS-LRM 的多图输入，Gamba 的技术重点是单图条件下如何高效、端到端地产生显式高斯资产；相对 DreamGaussian 一类优化式生成，它把耗时的逐实例优化替换为一次网络前向。\textbf{方法信息：} 代码：\url{https://github.com/skyworkai/gamba}；仓库提供权重和推理脚本，但训练教程仍标注为后续补充；GitHub 约 164 stars，Semantic Scholar 约 76 次引用。



% METHOD: 2412.12906
\paragraph{CATSplat（文本与空间双指导单图重建）（2024.12）}
CATSplat~\cite{2412_12906} 用双重先验补单图信息不足。它先用 UniDepth 得到深度图作为高斯中心候选，再由 ResNet 和多分辨率 Transformer 预测高斯参数。

文本指导来自视觉语言模型生成的场景描述，经 CLIP 文本嵌入通过交叉注意力注入图像特征；空间指导来自深度构建的 3D 点云特征，增强几何理解。相对 studentSplat 的教师伪标注，CATSplat 用文本上下文和深度点云先验。
ResNet 编码器提取的是深度条件化层次特征，而多分辨率 Transformer 则分别接收两类外部提示：文本提示提供对象和场景关系，空间提示来自深度反投影得到的 3D 点云特征。交叉注意力把这两类提示注入图像特征后，解码器再预测高斯位置、协方差、不透明度和颜色。它的核心增量不是新的渲染器，而是在单图高斯预测管线中加入语义上下文和几何候选，使网络在背面和遮挡区域有更强先验。
\textbf{方法信息：} 代码：\url{https://github.com/kuai-lab/iccv25_catsplat}。
\textbf{设计动机与边界：} 相对原始 3DGS 以独立高斯和图像重建误差为主的设定，CATSplat 把问题落在以下可实现的改动上：整体流程：(1) 使用预训练单目深度估计模型（UniDepth）预测深度图作为Gaussian中心候选。(2) ResNet编码器提取深度条件化的层次化图像特征。(3) 多分辨率Transformer通过交叉注意力融入两种先验：文本指导——用视觉-语言模型生成场景描述，通过交叉注意力将CLIP文本嵌入融入图像特征，补充上下文信息；空间指导——从深度估计构建3D点云，提取3D点特征通过交叉注意力注入图像特征，增强几何理解。(4) ResNet解码器从增强特征预测Gaussian参数。核心创新是在传统Gaussian参数预测管线中引入双重语义/几何先验。其设计目的不是单独增加一个模块，而是用上述表示、监督或渲染约束消除原流程在目标任务中的具体失配。提出双重指导机制：文本指导（VLM文本嵌入通过交叉注意力提供场景上下文）和空间指导（3D点特征提供几何先验）。在单视角3D场景重建中达到SOTA性能。适用时还应保留论文给出的边界：依赖VLM生成文本描述的质量；单视角固有的遮挡区域重建仍有限。
% METHOD: 2601.11772
\paragraph{studentSplat（师生单视角场景级 3DGS）（2026.01）}
studentSplat~\cite{2601_11772} 面向场景级单视角。教师模型是多视角 3DGS，通过跨视角匹配和三角化为每个像素提供高斯中心伪标注；学生模型只看单图，用几何损失和光度损失学习预测 3DGS。

外推网络在新视角重建损失前补全缺失上下文，缓解单视角背面和视野外区域。相对 Gamba 的物体重建，studentSplat 更关注场景级尺度歧义和外推。
教师阶段把多视角特征匹配和三角化结果转成每个像素的高斯中心伪标注，其中尺度只需在相对意义上稳定。学生模型接受单图输入，几何损失要求预测中心接近教师伪标注并保持深度梯度一致，光度损失则通过新视角渲染约束颜色和不透明度。外推网络在计算目标视角损失前补全不可见上下文，避免学生只复制输入可见表面。相对 Gamba 的物体重建，studentSplat 更关注场景级尺度歧义和外推区域。
\textbf{实现细节：} studentSplat（2026.01） 解决的核心问题是尺度歧义——从单张图像无法确定场景的绝对尺度，而直接用新视角监督训练单目模型会导致深度预测在任意尺度间摇摆。其设计灵感来自 §相应公式 节中多视角方法（如 DepthSplat）已具备可靠的几何推理能力，但单目方法无法直接利用这些能力。教师-学生架构。studentSplat 包含三个组件：(1) 多视角教师模型——以 DINOv2 为骨干的多视角 3DGS 网络，从多张输入图像预测可靠深度图；(2) 单目学生模型——以同一 DINOv2 骨干但仅接收单张图像，骨干分支沿 $z$ 轴定位高斯中心，细化分支用 CNN 特征精炼预测并解码其他高斯参数；(3) 外推器网络——预测缺失的场景上下文（输入视野之外的区域）。训练期间，教师模型为学生提供逐像素几何监督（深度损失），解决尺度歧义；外推器补充被遮挡区域的上下文信息。渐进式微调。studentSplat 还展示了一种实用策略：先用单目模型从一张输入图像和虚拟小偏移位姿生成若干新视角，再将这些新视角作为多视角教师模型的输入进行进一步精炼，形成"单目预测$\to$多视角精炼"的级联管线。在 RealEstate10K 和 ACID 的外推设置下，studentSplat 首次使单目方法的场景级新视角重建质量接近多视角方法。
\textbf{技术实现补充：} studentSplat（2026.01） 解决的核心问题是尺度歧义——从单张图像无法确定场景的绝对尺度，而直接用新视角监督训练单目模型会导致深度预测在任意尺度间摇摆。其设计灵感来自 §相应公式 节中多视角方法（如 DepthSplat）已具备可靠的几何推理能力，但单目方法无法直接利用这些能力。教师-学生架构。studentSplat 包含三个组件：(1) 多视角教师模型——以 DINOv2 为骨干的多视角 3DGS 网络，从多张输入图像预测可靠深度图；(2) 单目学生模型——以同一 DINOv2 骨干但仅接收单张图像，骨干分支沿 $z$ 轴定位高斯中心，细化分支用 CNN 特征精炼预测并解码其他高斯参数；(3) 外推器网络——预测缺失的场景上下文（输入视野之外的区域）。训练期间，教师模型为学生提供逐像素几何监督（深度损失），解决尺度歧义；外推器补充被遮挡区域的上下文信息。渐进式微调。studentSplat 还展示了一种实用策略：先用单目模型从一张输入图像和虚拟小偏移位姿生成若干新视角，再将这些新视角作为多视角教师模型的输入进行进一步精炼，形成"单目预测$\to$多视角精炼"的级联管线。在 RealEstate10K 和 ACID 的外推设置下，studentSplat 首次使单目方法的场景级新视角重建质量接近多视角方法。
\textbf{方法信息：} Amazon 参与了这篇论文。
\textbf{设计动机与边界：} 相对原始 3DGS 以独立高斯和图像重建误差为主的设定，studentSplat 把问题落在以下可实现的改动上：studentSplat采用师生架构：(1)教师模型：多视角3DGS模型，通过跨视角特征匹配和三角化估计每个像素的高斯中心（带隐式相对尺度），将数据集转换为带伪3D标注的版本；(2)学生模型：单视角3DGS模型，通过几何损失（教师高斯中心的L1损失+梯度一致性）和光度损失联合训练；(3)外推网络：在计算新视角重建损失前，补全缺失上下文，解决单视角模型的外推问题。还利用深度梯度正则化局部结构一致性。其设计目的不是单独增加一个模块，而是用上述表示、监督或渲染约束消除原流程在目标任务中的具体失配。提出studentSplat，首个场景级单视角3DGS方法。创新包括：(1)师生架构，多视角教师模型提供几何监督解决尺度模糊；(2)外推网络补全缺失场景上下文实现高质量外推。达到场景级单视角新视角重建的SOTA，同时作为自监督单视角深度估计方法展现竞争力。适用时还应保留论文给出的边界：重建质量仍受单视角固有模糊性限制；教师模型的质量影响学生模型性能；外推网络可能引入幻觉。

\section{领域化、紧凑与动态补充方法}

本节收纳几个不完全属于主流多视角前馈的分支：物体多模态表示、人体中心双目重建、紧凑 primitive 分配、未标定稀疏结构化重建和动态场景 densification。它们的共同点是把前馈 3DGS 从“通用新视角合成”推向更具体的部署约束。
% METHOD: 2312.12337
\paragraph{pixelSplat（前馈与可泛化重建）（2023.12）}
pixelSplat~\cite{2312_12337} 关注的问题是：用前馈或可泛化模型预测 3D Gaussian 场景 pixelSplat 是宽召回补全中发现的论文，主要根据标题、摘要和 arXiv 类别判定其与 3DGS 的关系。
相对原始 3DGS，pixelSplat 的变化不是只替换损失函数，而是把表示、初始化、约束或网络预测路径重新组织到任务流程中：该工作被召回是因为论文文本包含 Gaussian/3D Gaussian 相关表述；当前分类先确定其任务域和与 3DGS 的相关性，后续证据流水线再补代码、企业和引用信息。
\textbf{实现细节：} 具体实现可以按输入、核心模块和渲染/优化路径来读。输入端首先取得任务所需的图像、相机、文本、时间序列、语义或几何线索，并把这些线索转换成可以约束高斯中心、尺度、旋转、不透明度、颜色或附加特征的条件。核心结构中，方法使用的主要机制是：从摘要层面看，pixelSplat 围绕“用前馈或可泛化模型预测 3D Gaussian 场景”展开。若其正文或代码后续可读，应进一步核对它具体改动的是表示、初始化、渲染、优化、损失、动态建模、压缩或应用接口中的哪一环。 训练或推理时，这些模块通常先给出初始化或中间特征，再通过可微 Gaussian rasterization、重建损失、感知/语义约束、几何正则或任务损失回传到高斯参数；如果方法包含前馈网络，则网络输出的是高斯属性、变形场、压缩码、语义特征或控制变量，而不是只输出最终 2D 图像。渲染路径仍以显式高斯投影、深度排序和 alpha 合成为基础，只是在采样、属性预测、局部编辑、动态变形、压缩编码或下游任务接口处加入额外约束。这样写入本章时，它与前后方法的关系是：相对原始 3DGS 和本章已有锚点方法，它可以理解为 3DGS extends\_or\_uses 宽召回分类阶段的关系判断，后续需要按论文方法部分复核。
\textbf{证据与边界：} 本轮召回分类尚未逐项核对实验表格；后续证据补全阶段需要读取论文源码/PDF或项目页确认数据集、指标和对比方法。 该条来自宽召回批次；源码下载受 arXiv e-print 网络限制，部分论文暂时只能依据标题、摘要和类别分类，需要在入可信书前做人工/LLM精读复核。 本书只采用本地 paper pack、数据库分类、源码/引用登记和论文摘要中可核验的信息；没有本地证据支撑的速度、排行榜或开源完整性不在正文中强化。\textbf{方法信息：} 代码：\url{https://github.com/dcharatan/pixelsplat}；引用信号：732 citations (23.9/month, domain-relative cutoff 4.8/month, semantic\_scholar)；代码关注：dcharatan/pixelsplat has 1260 stars (41.0/month, domain-relative cutoff 32.0/month)。
\textbf{设计动机与边界：} 相对原始 3DGS 以独立高斯和图像重建误差为主的设定，pixelSplat 把问题落在以下可实现的改动上：从摘要层面看，pixelSplat 围绕“用前馈或可泛化模型预测 3D Gaussian 场景”展开。若其正文或代码后续可读，应进一步核对它具体改动的是表示、初始化、渲染、优化、损失、动态建模、压缩或应用接口中的哪一环。其设计目的不是单独增加一个模块，而是用上述表示、监督或渲染约束消除原流程在目标任务中的具体失配。该工作被召回是因为论文文本包含 Gaussian/3D Gaussian 相关表述；当前分类先确定其任务域和与 3DGS 的相关性，后续证据流水线再补代码、企业和引用信息。适用时还应保留论文给出的边界：该条来自宽召回批次；源码下载受 arXiv e-print 网络限制，部分论文暂时只能依据标题、摘要和类别分类，需要在入可信书前做人工/LLM精读复核。
% METHOD: 2410.18979
\paragraph{PixelGaussian（前馈与可泛化重建）（2024.10）}
PixelGaussian~\cite{2410_18979} 关注的问题是：用前馈或可泛化模型预测 3D Gaussian 场景 PixelGaussian 是宽召回补全中发现的论文，主要根据标题、摘要和 arXiv 类别判定其与 3DGS 的关系。
相对原始 3DGS，PixelGaussian 的变化不是只替换损失函数，而是把表示、初始化、约束或网络预测路径重新组织到任务流程中：该工作被召回是因为论文文本包含 Gaussian/3D Gaussian 相关表述；当前分类先确定其任务域和与 3DGS 的相关性，后续证据流水线再补代码、企业和引用信息。
\textbf{实现细节：} 具体实现可以按输入、核心模块和渲染/优化路径来读。输入端首先取得任务所需的图像、相机、文本、时间序列、语义或几何线索，并把这些线索转换成可以约束高斯中心、尺度、旋转、不透明度、颜色或附加特征的条件。核心结构中，方法使用的主要机制是：从摘要层面看，PixelGaussian 围绕“用前馈或可泛化模型预测 3D Gaussian 场景”展开。若其正文或代码后续可读，应进一步核对它具体改动的是表示、初始化、渲染、优化、损失、动态建模、压缩或应用接口中的哪一环。 训练或推理时，这些模块通常先给出初始化或中间特征，再通过可微 Gaussian rasterization、重建损失、感知/语义约束、几何正则或任务损失回传到高斯参数；如果方法包含前馈网络，则网络输出的是高斯属性、变形场、压缩码、语义特征或控制变量，而不是只输出最终 2D 图像。渲染路径仍以显式高斯投影、深度排序和 alpha 合成为基础，只是在采样、属性预测、局部编辑、动态变形、压缩编码或下游任务接口处加入额外约束。这样写入本章时，它与前后方法的关系是：相对原始 3DGS 和本章已有锚点方法，它可以理解为 3DGS extends\_or\_uses 宽召回分类阶段的关系判断，后续需要按论文方法部分复核。
\textbf{证据与边界：} 本轮召回分类尚未逐项核对实验表格；后续证据补全阶段需要读取论文源码/PDF或项目页确认数据集、指标和对比方法。 该条来自宽召回批次；源码下载受 arXiv e-print 网络限制，部分论文暂时只能依据标题、摘要和类别分类，需要在入可信书前做人工/LLM精读复核。 本书只采用本地 paper pack、数据库分类、源码/引用登记和论文摘要中可核验的信息；没有本地证据支撑的速度、排行榜或开源完整性不在正文中强化。\textbf{方法信息：} 代码：\url{https://github.com/barrybarry-smith/pixelgaussian}；代码关注：barrybarry-smith/pixelgaussian has 117 stars (5.7/month, domain-relative cutoff 32.0/month)。
\textbf{设计动机与边界：} 相对原始 3DGS 以独立高斯和图像重建误差为主的设定，PixelGaussian 把问题落在以下可实现的改动上：从摘要层面看，PixelGaussian 围绕“用前馈或可泛化模型预测 3D Gaussian 场景”展开。若其正文或代码后续可读，应进一步核对它具体改动的是表示、初始化、渲染、优化、损失、动态建模、压缩或应用接口中的哪一环。其设计目的不是单独增加一个模块，而是用上述表示、监督或渲染约束消除原流程在目标任务中的具体失配。该工作被召回是因为论文文本包含 Gaussian/3D Gaussian 相关表述；当前分类先确定其任务域和与 3DGS 的相关性，后续证据流水线再补代码、企业和引用信息。适用时还应保留论文给出的边界：该条来自宽召回批次；源码下载受 arXiv e-print 网络限制，部分论文暂时只能依据标题、摘要和类别分类，需要在入可信书前做人工/LLM精读复核。
% METHOD: 2506.04789
\paragraph{Object-X（多模态物体中心表示）（2025.06）}
Object-X~\cite{2506_04789} 学习物体中心的多模态 3D 表示。图像、点云和文本先被几何接地到 3D 体素网格，再与类别或描述属性融合，编码成紧凑 latent embedding。

训练上，自编码器用掩码 MSE 等目标学习能覆盖形状、外观和语义属性的紧凑嵌入；解码器再从该嵌入预测 3DGS 参数，实现高保真物体级渲染。同一个嵌入还能用于场景对齐、单图重建和定位。相对 GS-LRM，它的目标不是任意场景 NVS，而是把多模态物体资产压到一个可解码、可检索、可对齐的表示中。
\textbf{实现细节：} Object-X（2025.06） 提出多模态三维物体表示框架，将图像、点云、文本等多模态信息在三维体素网格中几何锚定。其核心设计是将每个物体编码为结构化潜在表示（Structured Latent Representation, SLat）：先将物体体素化为规范 3D 网格，在体素内聚合局部图像特征和点云几何特征，再通过稀疏 Transformer 编码器压缩为紧凑嵌入。该嵌入具有双重能力：(1) 通过 U-3DGS 解码器解码为三维高斯进行物体重建；(2) 直接作为描述符用于视觉定位、场景对齐和单图重建等下游任务，无需任务特定微调。
\textbf{设计动机与边界：} 相对原始 3DGS 以独立高斯和图像重建误差为主的设定，Object-X 把问题落在以下可实现的改动上：Object-X框架：(1)编码：将多模态输入（图像、点云、文本）几何接地到3D体素网格结构，与物体属性（类别标签、描述）融合，通过自编码器训练学习紧凑潜在嵌入；(2)解码：设计解码器从嵌入预测3DGS参数，实现高保真物体级渲染和重建。使用掩码MSE损失训练。嵌入同时可用于场景对齐、单图3D重建和定位。其设计目的不是单独增加一个模块，而是用上述表示、监督或渲染约束消除原流程在目标任务中的具体失配。提出Object-X框架，将多模态信息（图像、点云、文本）几何接地到3D体素网格，学习紧凑嵌入，可解码为高保真3DGS重建，同时支持场景对齐、单图重建和定位等下游任务。存储需求比传统方法少3-4个数量级。适用时还应保留论文给出的边界：对复杂拓扑物体的重建质量可能受限；编码过程需要多视角输入；下游任务性能依赖于嵌入质量。
\textbf{方法信息：} Microsoft 参与了这篇论文；代码：\url{https://github.com/gaiadilorenzo/object-x}。



% METHOD: 2511.22704
\paragraph{Splat-SAP（人体中心尺度感知点图）（2025.11）}
Splat-SAP~\cite{2511_22704} 面向双目人体中心场景。它先用 MASt3R 预测双视角 point maps，再通过迭代亲和力学习把规范空间点图转换到真实度量空间，解决 DUSt3R/MASt3R 归一化导致的尺度不稳定。

第二阶段将双视角点图投影到目标视角并做几何精化，再锚定高斯基元渲染。参考点 Plucker 坐标替代原始射线，以更好捕捉相对深度。相对通用 Stereo-GS，它更针对人体中心双目捕获。
尺度恢复依赖焦距和双目基线：这些相机内参被编码后，经自注意力和交叉注意力估计缩放因子；平移则由轻量卷积编码器和类似 RAFT 的迭代更新得到像素级偏移。整个几何阶段自监督训练，不需要真实 3D 标注。渲染阶段把精化点图投影到目标视角，通过立体匹配继续修正几何，再锚定深度、尺度、旋转、透明度和颜色等高斯属性。相对通用 Stereo-GS，它把尺度恢复、双目基线和人体中心场景作为主要约束。
\textbf{实现细节：} Splat-SAP（2025.11） 利用像素级 point map 重建表示几何，通过两阶段学习策略解决大稀疏双目输入的挑战。第一阶段基于 MASt3R 的 ViT 骨干提取双目图像 token，通过迭代亲和力学习将尺度不变的 point map 转换到真实尺度空间——这一步骤无需 3D 监督，仅通过渲染损失自监督训练。第二阶段将初始 point map 投影到目标视角形成高斯平面锚点，通过立体匹配约束精炼几何，在精炼后的几何上构建高斯基元进行渲染。整个粗到细管线仅用渲染损失训练，泛化到工业相机、GoPro 和手机三种不同相机类型。
\textbf{方法信息：} Alibaba 参与了这篇论文。
\textbf{设计动机与边界：} 相对原始 3DGS 以独立高斯和图像重建误差为主的设定，Splat-SAP 把问题落在以下可实现的改动上：Splat-SAP采用两阶段流程：(1)阶段1-尺度感知几何重建：基于MASt3R预测双视角点图，通过迭代亲和力学习将点图从规范空间转换到真实度量空间。缩放因子通过相机内参（焦距f）和基线距离d的编码特征经自注意力和交叉注意力计算；平移通过轻量卷积编码器和迭代更新（类似RAFT）计算像素级偏移。整个过程自监督训练，无需3D标注。(2)阶段2-几何精化与渲染：将精化后的双视角点图投影到目标视角，通过立体匹配精化几何，然后锚定高斯基元（深度、缩放、旋转、透明度、颜色）进行高质量图像渲染。采用参考点Plucker坐标替代原始Plucker射线以更好地捕捉相对深度和几何。其设计目的不是单独增加一个模块，而是用上述表示、监督或渲染约束消除原流程在目标任务中的具体失配。提出Splat-SAP，利用尺度感知点图重建几何表示，实现从稀疏双目输入的前馈高斯泼溅渲染。两阶段策略：阶段1通过迭代亲和力学习将点图从规范空间转换到真实度量空间（自监督，无需3D监督）；阶段2将双视角点图投影到目标视角并精化几何，锚定高斯基元实现高质量渲染。解决DUSt3R归一化导致的尺度不稳定问题。适用时还应保留论文给出的边界：依赖MASt3R的几何预测质量；人体运动可能导致规范空间中的相对深度差异引起抖动；双目输入限制了重建的完整性。
% METHOD: 2606.05102
\paragraph{ZipSplat（场景复杂度驱动高斯预算）（2026.06）}
\textbf{锚点细化：} ZipSplat：在 feed-forward 3DGS 的“紧凑前馈高斯预测”基础上推进，主要改变是：ZipSplat：token-based feed-forward模型，用多视图视觉token聚类生成紧凑scene tokens并解码为高斯组。它要解决的问题是：提出解耦Gaussian placement和像素网格的前馈模型，通过推理时聚类控制质量-速度折衷，减少平坦区域过量高斯。核心思想是：ZipSplat用scene token替代一像素一高斯的前馈范式，以更少高斯表达复杂场景。
ZipSplat~\cite{2606_05102} 是本章 ``像素网格与高斯预算解耦'' 的高关注新锚点候选。AnySplat、NoPoSplat 和 DepthSplat 这类方法虽然解决了相机或深度问题，但仍常从输入像素出发生成高斯：一面平墙和一簇细树枝只要占据相同像素面积，就会分到相近数量的 primitive。ZipSplat 的改变是在高斯生成前先压缩表示，把场景看作一组 scene tokens，而不是一张二维像素网格。

流程上，模型用多视图 foundation backbone 从稀疏、无位姿、无内参输入图像中提取 dense visual tokens。随后在特征空间做 k-means，把大量视觉 token 聚合成较少的 scene tokens；这些 tokens 经过 cross-attention 和 self-attention 继续融合跨视角信息，再由轻量 MLP 解码为若干组 3D Gaussians。由于高斯中心是 unconstrained 3D positions，而不是必须沿某个输入像素射线放置，ZipSplat 可以把容量集中到几何复杂或纹理丰富区域，同时合并多视图对同一表面的重复观测。训练时，为了避免自由 3D 位置失去几何锚定，方法加入几何监督，把解码出的高斯拉回有效场景表面。

一个关键工程点是压缩比例放在推理时控制：同一个训练好的模型可以通过调 k-means 保留 token 数，得到高质量重建或更小高斯集合，不必重新训练。论文报告在 DL3DV 和 RealEstate10K 上以约 $6\times$ 更少 Gaussians 超过 pose-free 基线，并在 Mip-NeRF 360、ScanNet++ 上做零样本验证。相对 SplatWeaver 的像素级基数专家，ZipSplat 的粒度更靠前，先把视觉证据压成 scene tokens 再决定高斯放置。\textbf{方法信息：} Microsoft 参与了这篇论文；代码：\url{https://github.com/cvg/ZipSplat}；仓库公开推理、评测、训练代码和预训练权重；GitHub 约 132 stars。

\textbf{实现细节：} ZipSplat（2026.06） 是紧凑前馈高斯预测路线中的新增代表。相较于 feed-forward 3DGS，主要改变是：ZipSplat：token-based feed-forward模型，用多视图视觉token聚类生成紧凑scene tokens并解码为高斯组。3DGS 的显式高斯表示和可微渲染能力保持为基础。其设计动机是：提出解耦Gaussian placement和像素网格的前馈模型，通过推理时聚类控制质量-速度折衷，减少平坦区域过量高斯。具体而言，多视图backbone提取dense tokens，经k-means聚合为紧凑scene tokens，再用cross/self-attention细化并由MLP解码成任意3D位置的高斯组。这给后续方法留下的比较点是：用token聚类生成紧凑高斯集合。边界在于：聚类质量影响高斯布局；复杂细节可能需要较高token预算；推理时聚类增加额外步骤。
\textbf{实现补充：} 实现上还要核对输入、输出、参数、特征、初始化、训练损失、正则、渲染或推理接口，以及优化中的剪枝、分裂、采样、编码或解码步骤；这些抓手决定该方法到底改了 3DGS 管线的哪一环。
% METHOD: 2606.28321
\paragraph{StructSplat（未标定稀疏结构化高斯）（2026.06）}
StructSplat~\cite{2606_28321} 也处理未标定稀疏视图，但它对 AnySplat 的改动方向不同。AnySplat 把相机、深度和高斯参数放进统一几何 Transformer，再通过体素化减少冗余；StructSplat 认为统一 backbone 容易把几何、语义和纹理线索缠在一起，导致未知相机输入下外观细节和几何稳定性相互牵制。它的核心是把这些线索组织成结构化表示，再分别用于几何定位和纹理恢复。

具体而言，StructSplat 在前馈网络中显式区分 geometry、semantic 和 texture cues。几何部分负责建立可渲染高斯的空间结构；semantic-aware 几何组件用语义先验帮助弱纹理或重复纹理区域保持全局一致；texture 分支则通过 pixel-aligned feature injection 把输入图像中的局部纹理重新注入到高斯外观预测中，避免高层全局特征抹掉细节。对于未知相机参数，模型还采用 camera alignment 策略，使不同视角的结构特征能在共同参考下融合，并降低相机信息泄漏到外观分支的风险。

因此，StructSplat 相对 AnySplat 的价值不在于再提出一个更大的相机头，而在于把未标定前馈 3DGS 拆成 ``结构先稳住、纹理再补齐'' 的流程。论文报告在 DL3DV 上比 AnySplat 有明显 PSNR 提升，并在 ACID、RealEstate10K 等跨数据集设置中验证泛化。\textbf{方法信息：} 代码：\url{https://github.com/J-C-Zhao/StructSplat}；仓库是 ECCV 2026 官方代码库，提供安装、数据和模型说明；GitHub 约 32 stars。
\textbf{设计动机与边界：} 相对原始 3DGS 以独立高斯和图像重建误差为主的设定，StructSplat 把问题落在以下可实现的改动上：通过 pixel-aligned feature injection 改善纹理建模，并引入 semantic-aware 几何组件，在未知相机参数输入下前馈预测可渲染 3D Gaussians。其设计目的不是单独增加一个模块，而是用上述表示、监督或渲染约束消除原流程在目标任务中的具体失配。提出组织 geometry、semantic 和 texture cues 的结构化表示，降低统一 backbone 中几何与外观纠缠。适用时还应保留论文给出的边界：无标定设定难度高，重建质量可能受相机估计、语义先验和输入视角覆盖限制。
% METHOD: 2606.29374
\paragraph{L2D2-GS（动态场景学习式致密化）（2026.06）}
L2D2-GS~\cite{2606_29374} 把前馈 3DGS 带到动态城市场景和自动驾驶仿真。静态前馈模型常把重建写成一次性参数回归：输入多视图图像，网络直接输出所有 Gaussian primitives。这个思路在动态街景里有两个瓶颈：高分辨率输入会让像素绑定高斯数量和显存同时膨胀；密集多视图观测包含运动物体、遮挡和时间变化，一次性回归很难决定哪里该增加 primitive、哪里该保持紧凑。

L2D2-GS 的改动是把 generalizable reconstruction 改写成可学习的迭代优化和 densification。框架采用 scene graph 表示，把城市环境拆成静态节点和动态节点，并用预定义 SE(3) 变换描述动态对象。每个节点由 voxel-wise neural Gaussians 参数化，重建模块用带稀疏卷积的 3D U-Net 反复优化节点内高斯属性。穿插在连续优化之间的是 policy-based densification：策略网络通过 attention 找出几何表达不足的区域，为候选位置分配新增高斯概率。

这个策略没有人工标注的 ``哪里应当致密化''。论文用可微渲染器把全局重建误差下降回溯到单个新增高斯，作为自监督 reward 训练 densification policy；同时用几何正则和 reparameterization 约束高斯属性的优化流形，避免早期 primitive 拉成长针状或被天空等低频区域错误补偿。相对 ReSplat 的渲染误差反馈，L2D2-GS 更强调动态场景中的 primitive 生成决策；相对 ZipSplat 的静态 token 压缩，它学习的是时间和运动条件下的 ``何处增密''。\textbf{方法信息：} Xiaomi 参与了这篇论文；本轮未发现 L2D2-GS 官方完整代码，DepthSplat 和 AnySplat 是相关工作或复用基础仓库，不作为本方法代码。


可泛化前馈 3DGS 的演进可以看作三次解耦。第一次是把逐场景优化解耦为网络前馈预测，MVSGaussian、GS-LRM、FreeSplat 和 DepthSplat 是这一层的锚点；第二次是把像素网格和高斯分布解耦，VolSplat、AnchorSplat、TokenGS、C3G、Off The Grid、SplatWeaver 等方法让 primitive 数量和位置更由 3D 场景复杂度决定；第三次是把相机标定假设解耦，PF3plat、NoPoSplat、FreeSplatter、AnySplat、YoNoSplat 等方法把位姿估计、内参估计和高斯预测放进同一前馈框架。读这章时不要只问模型多大，而要问它的高斯来自哪里：像素、深度、体素、锚点、token、全局场景记忆，还是动态 densification。

\textbf{设计动机与边界：} 相对原始 3DGS 以独立高斯和图像重建误差为主的设定，L2D2-GS 把问题落在以下可实现的改动上：在多视图动态输入下，逐步分配和细化 Gaussian primitives，使密集观测融合更稳健，兼顾高保真重建和比逐场优化更快的推理。其设计目的不是单独增加一个模块，而是用上述表示、监督或渲染约束消除原流程在目标任务中的具体失配。将 feedforward Gaussian reconstruction 从一次性参数回归转为可学习的多步 densification，提高动态城市场景的可扩展性。适用时还应保留论文给出的边界：迭代致密化仍可能增加推理步骤；动态城市长尾对象和遮挡场景需要更多验证。
\textbf{尾项实现澄清：} L2D2-GS 针对一次前向回归在高分辨率动态城市场景中内存重、密集观测融合不稳的问题。它把前馈高斯重建改为多步优化和可学习致密化：全局重建增益形成自监督奖励，指导局部 primitive 何时增密；几何正则和重参数化限制优化流形，减少早期伪影固化。相对一次性 feedforward 锚点，它在推理中增加受控迭代，以较少高斯换取更稳定的动态多视图融合。
